% ============================================================
%  Géolocalisation des Eglises EEC — ENSPY 2025-2026
%  pdfLaTeX + biber
%  Compilation : pdflatex → biber → pdflatex → pdflatex
% ============================================================
\documentclass[12pt, a4paper]{report}

\usepackage[utf8]{inputenc}
\usepackage[T1]{fontenc}
\usepackage[french]{babel}
\usepackage{csquotes}
\usepackage[a4paper, top=2.5cm, bottom=2.5cm,
            left=3cm, right=2.5cm, headheight=14pt]{geometry}
\usepackage{lmodern}
\usepackage{microtype}
\usepackage{setspace}
\setlength{\parindent}{1.2cm}
\setlength{\parskip}{0pt}

% ---- Couleurs ----
\usepackage{xcolor}
\definecolor{eecgreen}{RGB}{27,94,32}
\definecolor{eccdarkgreen}{RGB}{17,60,20}
\definecolor{eecyellow}{RGB}{249,168,37}
\definecolor{eeclightgreen}{RGB}{232,245,233}
\definecolor{bordergold}{RGB}{212,160,23}

% ---- Bordure double sur chaque page ----
\usepackage{eso-pic}
\usepackage{tikz}
\AddToShipoutPictureBG{%
  \begin{tikzpicture}[remember picture, overlay]
    \draw[color=bordergold, line width=2pt]
      ([xshift=10pt,yshift=-10pt]current page.north west)
      rectangle
      ([xshift=-10pt,yshift=10pt]current page.south east);
    \draw[color=eecgreen, line width=0.7pt]
      ([xshift=15pt,yshift=-15pt]current page.north west)
      rectangle
      ([xshift=-15pt,yshift=15pt]current page.south east);
  \end{tikzpicture}}

% ---- En-têtes et pieds — texte noir ----
\usepackage{fancyhdr}
\pagestyle{fancy}
\fancyhf{}
\fancyhead[L]{\small\textit{Géolocalisation des Eglises EEC --- ENSPY 2025-2026}}
\fancyhead[R]{\small\thepage}
\fancyfoot[C]{\small\textit{Département de Génie Informatique}}
\renewcommand{\headrulewidth}{0.4pt}
\renewcommand{\footrulewidth}{0pt}

% ---- Titres de chapitre ----
\usepackage{titlesec}
% Chapitres numérotés — boîte décorative verte + numéro + titre en petites capitales
\titleformat{\chapter}[display]{%
  \normalfont%
}{%
  \begin{tikzpicture}[remember picture,overlay]%
    \node[rotate=90, anchor=center,
          text=eecgreen!80!black,
          font=\footnotesize\bfseries\scshape,
          inner sep=0pt]
      at ([xshift=-72pt, yshift=-55pt]current page.north east)
      {Chapitre};%
    \fill[eecgreen]
      ([xshift=-60pt, yshift=-10pt]current page.north east)
      rectangle
      ([xshift=0pt, yshift=-100pt]current page.north east);%
    \node[text=white, font=\fontsize{68}{68}\selectfont\bfseries, anchor=center]
      at ([xshift=-30pt, yshift=-55pt]current page.north east)
      {\thechapter};%
  \end{tikzpicture}%
}{10pt}{%
  \centering\normalfont\large\bfseries\scshape\color{eecgreen}%
}[\vspace{2pt}{\color{eecyellow}\hrule height 2.5pt}\vspace{4pt}]
\titlespacing{\chapter}{0pt}{26pt}{14pt}

% Chapitres non numérotés (\chapter*) — style simple sans décoration
\titleformat{name=\chapter,numberless}[block]{%
  \normalfont\Large\bfseries\color{eecgreen}%
}{}{0pt}{\MakeUppercase}[%
  \vspace{3pt}{\color{eecyellow}\hrule height 2.5pt}\vspace{5pt}%
]
\titlespacing{name=\chapter,numberless}{0pt}{0pt}{10pt}

% Sections et sous-sections — hiérarchie verte grasse
\titleformat{\section}
  {\normalfont\large\bfseries\color{eccdarkgreen}}
  {\thesection.}{8pt}{}
  [\vspace{1pt}{\color{eecgreen!40!white}\hrule height 0.6pt}\vspace{2pt}]
\titleformat{\subsection}
  {\normalfont\normalsize\bfseries\color{eecgreen}}
  {\thesubsection.}{7pt}{}
\titleformat{\subsubsection}
  {\normalfont\normalsize\bfseries\color{eecgreen!90!black}}
  {\thesubsubsection.}{6pt}{}

% ---- Tableaux ----
\usepackage{array}
\usepackage{colortbl}
\usepackage{booktabs}
\usepackage{longtable}

% ---- Figures et flottants ----
\usepackage{graphicx}
\usepackage{float}

% ---- Listes ----
\usepackage{enumitem}

% ---- Lettrine (lettre ornée de début de chapitre) ----
\usepackage{lettrine}
\renewcommand{\LettrineFontHook}{\color{eecgreen}\bfseries}
\setcounter{DefaultLines}{2}
\setlength{\DefaultFindent}{4pt}
\setlength{\DefaultNfindent}{0pt}

% ---- Symboles ----
\usepackage{pifont}
\usepackage{amssymb}

% ---- Liens ----
\usepackage[colorlinks=true, linkcolor=eecgreen,
            urlcolor=eecgreen, citecolor=eecgreen]{hyperref}

% ---- Bibliographie (BibLaTeX + biber) ----
\usepackage[style=numeric, sorting=none, backend=biber,
            bibencoding=utf8]{biblatex}
\addbibresource{references.bib}

% ---- Sommaire local par chapitre (minitoc — doit être APRÈS hyperref) ----
\usepackage{minitoc}
\mtcsettitle{minitoc}{\normalsize\bfseries\color{eecgreen}Sommaire}
\setcounter{minitocdepth}{3}

% ---- Symboles tableaux comparatifs ----
\newcommand{\oui}{\textcolor{eecgreen}{\ding{51}}}
\newcommand{\non}{\textcolor{red}{\ding{55}}}
\newcommand{\partiel}{\textcolor{orange}{$\sim$}}

% ---- Mots-clés techniques (premier emploi) ----
\newcommand{\kw}[1]{\textbf{#1}}

% ---- Plan de chapitre (boîte colorée en tête de chapitre) ----
\newenvironment{chapplan}{%
  \vspace{6pt}%
  \noindent\colorbox{eeclightgreen!70}{%
    \begin{minipage}[t]{\dimexpr\linewidth-2\fboxsep}%
      \vspace{5pt}%
      \noindent{\small\bfseries\color{eccdarkgreen}%
        \ding{228}\quad Plan du chapitre}%
      \vspace{3pt}\\%
      {\color{eecgreen!60!white}\hrule height 0.5pt}%
      \vspace{2pt}%
      \begin{itemize}[%
        leftmargin=1.4cm, itemsep=1pt, topsep=2pt, parsep=0pt,%
        label={\small\color{eecgreen}\ding{118}}]%
      \small%
}{%
      \end{itemize}%
      \vspace{4pt}%
    \end{minipage}%
  }%
  \vspace{8pt}%
}

% ============================================================
\begin{document}
\pagenumbering{roman}
\setstretch{1.3}

% ============================================================
%  RÉSUMÉ
% ============================================================
\chapter*{Résumé}
\addcontentsline{toc}{chapter}{Résumé}
\markboth{Résumé}{Résumé}

\noindent
L'Église Évangélique du Cameroun (EEC) regroupe 565~paroisses, 311~œuvres
sociales et éducatives, ainsi que 685~ouvriers ecclésiastiques répartis dans
22~Régions Synodales et 137~districts couvrant l'ensemble du territoire
national, sans que l'institution ne dispose d'aucun outil numérique centralisé
permettant leur visualisation géographique ni leur administration hiérarchisée.
Cette lacune compromet la capacité du Bureau National à piloter ses ressources
territoriales, à produire des statistiques consolidées et à assurer une mise à
jour contrôlée des données. Ce travail présente la conception et le développement
d'une plateforme web cartographique institutionnelle répondant à ces besoins.
L'approche adoptée s'appuie sur une architecture à services découplés~: Django~5
et Django REST Framework assurent le backend et les API, Next.js~14 constitue
l'interface web, PostgreSQL avec l'extension PostGIS gère la persistance des
données spatiales, GeoServer assure la publication des couches cartographiques
selon les standards OGC, et Leaflet.js permet la visualisation interactive côté
navigateur~; l'ensemble est conteneurisé avec Docker Compose. La plateforme
ainsi développée offre une carte interactive avec regroupement automatique des
marqueurs, des filtres hiérarchiques par région, district et type d'entité, un
panneau de détail par entité, ainsi qu'un module d'administration multi-niveaux
fondé sur un contrôle d'accès basé sur les rôles (RBAC) à quatre profils~:
national, régional, district et paroissial. Par rapport aux solutions existantes
analysées, ce système se distingue par son alignement strict avec le référentiel
hiérarchique synodal de l'EEC, l'utilisation de GeoServer pour la diffusion
normalisée des données géographiques, et un déploiement entièrement conteneurisé
garantissant la reproductibilité de l'environnement d'exécution.

\smallskip
\noindent\textbf{Mots-clés~:}
géolocalisation, système d'information géographique, EEC~Cameroun,
Django REST Framework, Leaflet.

\vspace{0.8cm}
\noindent{\color{eecgreen}\rule{\linewidth}{0.8pt}}
\vspace{0.3cm}

% ============================================================
%  ABSTRACT
% ============================================================
\chapter*{Abstract}
\addcontentsline{toc}{chapter}{Abstract}
\markboth{Abstract}{Abstract}

\noindent
The Evangelical Church of Cameroon (EEC) encompasses 565~parishes,
311~social and educational institutions, and 685~ordained workers spread across
22~Synodal Regions and 137~districts covering the entire national territory,
yet the institution lacks any centralized digital system for their geographic
visualization or hierarchical administration. This gap undermines the National
Bureau's ability to monitor its territorial resources, generate consolidated
statistics, and maintain controlled data updates. This work presents the design
and development of an institutional cartographic web platform addressing these
needs. The approach relies on a decoupled service architecture~: Django~5 and
Django REST Framework handle the backend and API layer, Next.js~14 provides the
web interface, PostgreSQL with the PostGIS extension manages spatial data
persistence, GeoServer handles OGC-compliant cartographic layer publication, and
Leaflet.js enables client-side interactive mapping~; the entire stack is
containerized using Docker Compose. The resulting platform delivers an
interactive map with automatic marker clustering, hierarchical filters by region,
district, and entity type, a per-entity detail panel, and a multi-level
administration module enforcing role-based access control (RBAC) across four
profiles~: national, regional, district, and parish. Compared to the existing
solutions reviewed in this work, the system is distinguished by its strict
alignment with the EEC synodal hierarchical framework, the use of GeoServer for
standards-compliant geographic data publication, and a fully containerized
deployment ensuring environment reproducibility.

\smallskip
\noindent\textbf{Keywords~:}
geolocation, geographic information system, EEC~Cameroon,
Django REST Framework, Leaflet.

\clearpage

% ============================================================
%  SIGLES ET ABRÉVIATIONS
% ============================================================
\chapter*{Sigles et Abréviations}
\addcontentsline{toc}{chapter}{Sigles et Abréviations}
\markboth{Sigles et Abréviations}{Sigles et Abréviations}

\noindent
Le tableau ci-dessous recense par ordre alphabétique les sigles et abréviations
employés dans ce document.

\vspace{4pt}

\begingroup
\setstretch{1.0}
\renewcommand{\arraystretch}{1.12}
\noindent
\begin{tabular}{%
  >{\bfseries}p{3.2cm}%
  p{10.8cm}}
  \rowcolor{eecgreen}
  \multicolumn{1}{>{\bfseries\color{white}}p{3.2cm}}{Sigle} &
  \multicolumn{1}{>{\bfseries\color{white}}p{10.8cm}}{Signification complète} \\
  \rowcolor{eeclightgreen!60}
  API     & Application Programming Interface \\
  CORS    & Cross-Origin Resource Sharing \\
  \rowcolor{eeclightgreen!60}
  CRUD    & Create, Read, Update, Delete \\
  CSS     & Cascading Style Sheets \\
  \rowcolor{eeclightgreen!60}
  DRF     & Django REST Framework \\
  EEC     & Église Évangélique du Cameroun \\
  \rowcolor{eeclightgreen!60}
  ENSPY   & École Nationale Supérieure Polytechnique de Yaoundé \\
  GIS     & Geographic Information System \\
  \rowcolor{eeclightgreen!60}
  GPS     & Global Positioning System (Système de Positionnement Global) \\
  HTML    & HyperText Markup Language \\
  \rowcolor{eeclightgreen!60}
  HTTP    & HyperText Transfer Protocol \\
  HTTPS   & HyperText Transfer Protocol Secure \\
  \rowcolor{eeclightgreen!60}
  JSON    & JavaScript Object Notation \\
  JWT     & JSON Web Token \\
  \rowcolor{eeclightgreen!60}
  OGC     & Open Geospatial Consortium \\
  ORM     & Object-Relational Mapping \\
  \rowcolor{eeclightgreen!60}
  PDF     & Portable Document Format \\
  PostGIS & PostgreSQL Spatial Extension \\
  \rowcolor{eeclightgreen!60}
  RBAC    & Role-Based Access Control (Contrôle d'accès basé sur les rôles) \\
  REST    & Representational State Transfer \\
  \rowcolor{eeclightgreen!60}
  SIG     & Système d'Information Géographique \\
  SLD     & Styled Layer Descriptor \\
  \rowcolor{eeclightgreen!60}
  SND30   & Stratégie Nationale de Développement 2020--2030 \\
  SQL     & Structured Query Language \\
  \rowcolor{eeclightgreen!60}
  SRS     & Spatial Reference System \\
  SSL     & Secure Sockets Layer \\
  \rowcolor{eeclightgreen!60}
  SSR     & Server-Side Rendering \\
  TOTP    & Time-based One-Time Password \\
  \rowcolor{eeclightgreen!60}
  UI      & User Interface (Interface utilisateur) \\
  URL     & Uniform Resource Locator \\
  \rowcolor{eeclightgreen!60}
  UX      & User Experience (Expérience utilisateur) \\
  WFS     & Web Feature Service \\
  \rowcolor{eeclightgreen!60}
  WGS84   & World Geodetic System 1984 \\
  WMS     & Web Map Service \\
  \rowcolor{eeclightgreen!60}
  WSL     & Windows Subsystem for Linux \\
\end{tabular}
\endgroup

\clearpage

% ============================================================
%  GLOSSAIRE
% ============================================================
\chapter*{Glossaire}
\addcontentsline{toc}{chapter}{Glossaire}
\markboth{Glossaire}{Glossaire}

\noindent
Ce glossaire définit les termes techniques spécialisés employés dans ce document.

\vspace{6pt}

\noindent\textbf{Cluster de marqueurs~:}
Technique de regroupement visuel sur une carte qui agrège les marqueurs
géographiquement proches en un seul indicateur numérique. Lorsque l'utilisateur
zoome, le cluster éclate pour révéler les marqueurs individuels. Cette technique
améliore la lisibilité d'une carte comportant un grand nombre de points.

\medskip
\noindent\textbf{Conteneurisation~:}
Approche de déploiement logiciel qui isole une application et toutes ses
dépendances dans une unité standardisée appelée conteneur. Dans ce projet, Docker
et Docker Compose orchestrent l'ensemble des services (backend, base de données,
serveur cartographique, cache, frontend) de manière reproductible et indépendante
de l'environnement hôte.

\medskip
\noindent\textbf{GeoServer~:}
Serveur cartographique open source conforme aux standards de l'Open Geospatial
Consortium (OGC). Il permet de publier, gérer et diffuser des données
géographiques sous forme de services WMS et WFS, accessibles depuis n'importe
quel client cartographique compatible.

\medskip
\noindent\textbf{Géolocalisation~:}
Processus de détermination et d'enregistrement de la position géographique d'une
entité à l'aide de coordonnées GPS (latitude, longitude). Dans ce projet, la
géolocalisation désigne l'association de chaque paroisse et œuvre de l'EEC à des
coordonnées précises sur le territoire camerounais.

\medskip
\noindent\textbf{PostGIS~:}
Extension spatiale pour le système de gestion de base de données relationnelle
PostgreSQL. PostGIS ajoute la prise en charge des types de données géographiques
(points, lignes, polygones) et des fonctions de calcul spatial (distances,
intersections, projections). Il constitue le standard de facto pour les bases de
données spatiales open source.

\medskip
\noindent\textbf{RBAC --- Contrôle d'accès basé sur les rôles~:}
Modèle de sécurité informatique dans lequel les droits d'accès aux ressources
d'un système sont déterminés par le rôle attribué à chaque utilisateur, et non
par son identité propre. Quatre rôles sont définis~: national, régional, district
et paroissial, chacun disposant d'un périmètre d'action délimité.

\medskip
\noindent\textbf{Région synodale~:}
Entité administrative de l'Église Évangélique du Cameroun regroupant un ensemble
de districts et de paroisses dans une zone géographique définie. L'EEC compte
22~régions synodales couvrant l'ensemble du territoire camerounais. La région
synodale est le premier niveau hiérarchique sous le Bureau National.

\medskip
\noindent\textbf{Shapefile~:}
Format de fichier géospatial vectoriel développé par ESRI, largement utilisé pour
la représentation de données géographiques. Un shapefile est constitué de plusieurs
fichiers complémentaires (\texttt{.shp}, \texttt{.dbf}, \texttt{.prj},
\texttt{.shx}) qui encodent la géométrie, les attributs et le système de
coordonnées des entités géographiques.

\medskip
\noindent\textbf{SLD --- Styled Layer Descriptor~:}
Standard OGC permettant de décrire la symbolisation et l'apparence d'une couche
géographique publiée via WMS. Les fichiers SLD définissent les couleurs, tailles
et règles d'affichage des entités géographiques sur la carte.

\medskip
\noindent\textbf{Tuile cartographique~:}
Image raster de dimensions fixes (généralement 256~$\times$~256~pixels)
représentant une portion de carte à un niveau de zoom donné. Les bibliothèques
cartographiques comme Leaflet assemblent dynamiquement ces tuiles pour former la
carte visible par l'utilisateur.

\medskip
\noindent\textbf{WFS --- Web Feature Service~:}
Standard OGC permettant de requêter et de récupérer des données géographiques
vectorielles (points, lignes, polygones) depuis un serveur cartographique,
généralement au format GeoJSON ou GML. Contrairement au WMS qui retourne des
images, le WFS retourne des données brutes manipulables côté client.

\medskip
\noindent\textbf{WMS --- Web Map Service~:}
Standard OGC permettant de requêter des cartes sous forme d'images raster depuis
un serveur cartographique. Le client envoie une requête HTTP spécifiant l'étendue
géographique et les couches souhaitées~; le serveur retourne une image PNG ou
JPEG correspondante, prête à être superposée sur la carte.

\clearpage

% ============================================================
%  TABLE DES MATIÈRES
% ============================================================
\dominitoc
\tableofcontents
\clearpage

% ============================================================
%  LISTE DES FIGURES
% ============================================================
\listoffigures
\addcontentsline{toc}{chapter}{Liste des figures}
\clearpage

% ============================================================
%  LISTE DES TABLEAUX
% ============================================================
\listoftables
\addcontentsline{toc}{chapter}{Liste des tableaux}
\clearpage

% ============================================================
%  CORPS DU RAPPORT — numérotation arabe
% ============================================================
\pagenumbering{arabic}
\setcounter{page}{1}

% ============================================================
%  INTRODUCTION GÉNÉRALE
% ============================================================
\chapter*{Introduction Générale}
\addcontentsline{toc}{chapter}{Introduction Générale}
\markboth{Introduction Générale}{Introduction Générale}

La transformation numérique des institutions publiques et religieuses constitue
l'un des axes stratégiques du développement durable en Afrique subsaharienne.
Selon l'Union Internationale des Télécommunications, le taux de pénétration
d'Internet en Afrique subsaharienne a atteint 36~\% en 2023~\cite{itu2023},
progressant de plus de dix points en cinq ans grâce au déploiement massif des
réseaux mobiles. La région comptait en 2023 plus de 615~millions d'abonnés à la
téléphonie mobile~\cite{gsma2023}, une dynamique qui confère aux technologies web
un rôle central dans la modernisation des organisations de toute nature. Au
Cameroun, cette tendance est amplifiée par la Stratégie Nationale de
Développement 2020--2030 (SND30), qui identifie explicitement la numérisation de
l'administration publique et des institutions nationales comme l'un des leviers
prioritaires du développement économique du pays~\cite{snd2020}.

Dans ce contexte, les Systèmes d'Information Géographique (SIG) se sont imposés
comme des outils incontournables pour la gestion et la visualisation des données
territoriales. Longley et al.~définissent un SIG comme un système informatique
conçu pour la collecte, le stockage, la manipulation, l'analyse et la
représentation de données géoréférencées~\cite{longley2015}. À l'échelle
mondiale, ces plateformes sont déployées dans des domaines aussi variés que la
santé publique, la gestion urbaine, la logistique humanitaire et la planification
administrative. La standardisation des échanges de données géographiques par
l'Open Geospatial Consortium (OGC) --- notamment les protocoles Web Map Service
(WMS)~\cite{ogcwms2006} et Web Feature Service (WFS)~\cite{ogcwfs2010} --- a
rendu ces technologies accessibles à des organisations de toute taille et de tout
secteur, y compris les institutions confessionnelles. En Afrique centrale, les
organisations religieuses constituent souvent le premier réseau d'établissements
scolaires, de centres de santé et de structures d'encadrement rural présent dans
les zones éloignées des grandes agglomérations~\cite{gifford2015}. Leur capacité
à gérer efficacement des ressources géographiquement dispersées à travers de
larges territoires nationaux constitue un enjeu direct pour le développement
des communautés qu'elles desservent.

L'Église Évangélique du Cameroun (EEC), fondée en 1957 dans la continuité des
missions protestantes bâloises et parisiennes, est l'une des plus grandes
institutions protestantes d'Afrique centrale. Son organisation territoriale
repose sur une hiérarchie administrative à quatre niveaux~: le Bureau National,
qui assure la direction centrale de l'institution~; les 22~Régions Synodales, qui
couvrent l'ensemble du territoire camerounais~; les 137~districts ecclésiaux, qui
coordonnent les activités locales~; et les 565~paroisses, qui constituent le
maillage de base de la présence de l'EEC sur le terrain. À ces entités s'ajoutent
311~œuvres sociales et éducatives --- établissements scolaires, centres médicaux,
exploitations agropastorales, immeubles et terrains --- et plus de
685~ouvriers ecclésiastiques (pasteurs, évangélistes, diacres permanents) répartis
sur l'ensemble de ce réseau.

Malgré cette envergure nationale considérable, l'EEC ne dispose à ce jour d'aucun
système numérique centralisé permettant la visualisation géographique de ses
entités, leur administration hiérarchisée selon les périmètres synodaux, ni la
consolidation automatique de leurs données statistiques. La localisation des
paroisses, la gestion des affectations des ouvriers et le suivi des statistiques
annuelles sont assurés manuellement via des fichiers tableur transmis
périodiquement entre niveaux hiérarchiques. Cette situation génère plusieurs
problèmes opérationnels concrets~: l'absence de vue géographique d'ensemble
empêche le Bureau National d'identifier les zones sous-dotées en ressources
pastorales ou d'optimiser l'affectation des ouvriers~; la dispersion des données
dans des fichiers non interconnectés rend toute consolidation statistique nationale
laborieuse et sujette aux erreurs~; l'absence de contrôle d'accès formalisé
compromet l'intégrité des informations institutionnelles~; et l'audit des données
sources révèle que 127~paroisses, soit 22~\% de l'effectif total, ne disposent
d'aucune coordonnée GPS valide, ce qui illustre l'inexistence d'un processus
structuré de collecte et de validation des données de localisation.

Plusieurs catégories de solutions pourraient théoriquement répondre à ces besoins.
Les plateformes SIG généralistes comme ArcGIS~Online~\cite{arcgisonline2024} ou
Google~Maps~Platform~\cite{googlemaps2024} offrent des capacités cartographiques
avancées mais demeurent des outils génériques, soumis à des licences commerciales
onéreuses et dépourvus de toute logique hiérarchique institutionnelle. Les outils
open~source comme QGIS~\cite{qgis2024} ou OpenLayers permettent la visualisation
cartographique avancée mais ne constituent pas des plateformes web multi-utilisateurs
dotées de contrôle d'accès et de gestion de données métier. Une première version
d'une plateforme de géolocalisation EEC avait été développée antérieurement --- elle
offrait une cartographie de base mais reposait sur des fichiers GeoJSON statiques,
ne disposait pas d'un mécanisme d'authentification par sessions sécurisées, et son
contrôle d'accès ne correspondait pas à la granularité de la hiérarchie synodale.
Aucune des solutions recensées ne répond simultanément aux exigences d'une
plateforme institutionnelle sécurisée, intégrée, alignée sur la structure synodale
de l'EEC et déployable dans un environnement à infrastructure informatique limitée.

\medskip

Face à ces constats, la problématique suivante guide ce travail~:

\vspace{6pt}
\noindent\colorbox{eeclightgreen!80}{\parbox{0.88\linewidth}{%
  \vspace{6pt}%
  \noindent\textit{\textbf{Question de recherche~:} Comment concevoir et
  implémenter une plateforme web cartographique institutionnelle répondant
  simultanément aux exigences de visualisation géographique interactive,
  d'administration hiérarchisée à quatre niveaux de l'autorité synodale et de
  sécurité des données de l'Église Évangélique du Cameroun~?}%
  \vspace{6pt}%
}}
\vspace{6pt}

Nous posons l'hypothèse qu'une architecture à services découplés --- combinant un
backend Django~5 avec GeoDjango et Django REST Framework pour la logique métier et
les interfaces de programmation (API), une base de données PostgreSQL avec
l'extension PostGIS pour la persistance des données spatiales, un serveur
cartographique GeoServer~2.26 pour la publication des couches géographiques selon
les standards OGC (WMS et WFS), et une interface Next.js~14 pour la consultation
publique et l'administration --- permettrait de répondre à cette problématique. La
sécurité reposerait sur un mécanisme d'authentification par sessions Django
(\textit{HttpOnly}, \textit{Secure}, \textit{SameSite}) et un contrôle d'accès
basé sur les rôles (RBAC) à quatre niveaux hiérarchiques aligné sur la structure
synodale. La conteneurisation de l'ensemble via Docker Compose garantirait la
reproductibilité de l'environnement d'exécution et la facilité du déploiement.

\medskip

Pour répondre à cette problématique, ce travail se fixe les objectifs suivants~:

\begin{enumerate}[leftmargin=1.5cm, itemsep=2pt, topsep=4pt]
  \item \textbf{Modéliser} la structure de données hiérarchique de l'EEC
    (régions synodales, districts, paroisses, œuvres sociales, ouvriers
    ecclésiastiques) dans un schéma de base de données relationnelle-spatiale
    PostGIS intégrant les contraintes d'intégrité référentielle et les
    informations de géolocalisation~;

  \item \textbf{Développer} un backend API REST sécurisé avec Django~5 et
    Django REST Framework, implémentant un contrôle d'accès basé sur les rôles
    à quatre niveaux --- national, régional, district et paroissial --- et un
    journal d'audit automatique de toutes les opérations d'écriture~;

  \item \textbf{Configurer} un serveur cartographique GeoServer~2.26 publiant
    les couches géographiques de l'EEC selon les standards OGC WMS~1.3.0 et
    WFS~2.0.0, avec des styles SLD aux couleurs institutionnelles de l'EEC~;

  \item \textbf{Implémenter} une interface cartographique interactive basée sur
    Leaflet.js avec regroupement automatique des marqueurs (\textit{clustering}),
    filtres hiérarchiques multicritères, popups de détail par entité et contrôle
    dynamique des couches cartographiques~;

  \item \textbf{Réaliser} un module d'administration multi-niveaux en Next.js~14
    permettant le CRUD des entités selon le périmètre du rôle, l'import des
    données depuis les fichiers sources, l'export PDF et Excel, et la
    visualisation des statistiques avancées par région et par année~;

  \item \textbf{Conteneuriser} l'ensemble de la plateforme avec Docker Compose
    en cinq services afin de garantir la reproductibilité de l'environnement
    d'exécution et de simplifier le déploiement en production.
\end{enumerate}

\medskip

La démarche adoptée s'organise en phases séquentielles et itératives. Une première
phase d'audit des données sources --- fichiers Excel paroisses, ouvriers et
œuvres, et Shapefile des régions synodales --- permet de valider la structure
réelle des données avant toute modélisation, notamment d'identifier les problèmes
de qualité~: inversion des coordonnées GPS dans les fichiers sources, incohérences
de nommage entre l'Excel et le Shapefile, paroisses sans géolocalisation. Une
deuxième phase de conception définit le schéma de base de données, l'architecture
des services et les modèles UML. L'implémentation procède ensuite couche par
couche~: modèles de données et scripts d'import, API REST avec sécurisation RBAC,
publication GeoServer des couches cartographiques avec styles SLD, développement
du frontend public et de l'interface d'administration. Une validation fonctionnelle
couvre les scénarios d'utilisation des quatre profils d'accès et les tests de
sécurité conformément aux recommandations OWASP~\cite{owasp2021}.

\medskip

La suite de ce rapport est structurée comme suit. Le \textbf{Chapitre~1} présente
les concepts généraux fondamentaux --- SIG, standards OGC, architecture REST,
contrôle d'accès basé sur les rôles --- et expose l'état de l'art des solutions
de gestion cartographique institutionnelle existantes~; il conclut par un tableau
comparatif et le positionnement de notre solution par rapport à l'existant. Le
\textbf{Chapitre~2} expose l'analyse des besoins à travers un tableau d'exigences
fonctionnelles et non-fonctionnelles ainsi que les diagrammes UML nécessaires à la
compréhension du système, puis présente la conception de l'architecture technique
retenue. Le \textbf{Chapitre~3} détaille l'environnement d'implémentation et les
choix technologiques, présente les résultats à travers les interfaces fonctionnelles
de la plateforme, et propose une discussion critique des forces et des limites de la
solution réalisée. Une \textbf{Conclusion Générale} clôt ce document en synthétisant
les apports du travail, ses limites et les perspectives d'amélioration et de
déploiement.

\clearpage

% ============================================================
%  CHAPITRE 1 — CONCEPTS GÉNÉRAUX ET ÉTAT DE L'ART
% ============================================================
\chapter{Concepts Généraux et État de l'Art}

\begin{chapplan}
  \item \textbf{Section~\ref{sec:concepts} — Concepts Généraux}~: Système
    d'Information Géographique, standards OGC (WMS\,/\,WFS), base de données
    spatiale PostGIS, architecture REST, contrôle d'accès RBAC et conteneurisation Docker.
  \item \textbf{Section~\ref{sec:etat-art} — État de l'Art}~: analyse comparative
    de quatre solutions (GeoEEC~v1.0, ArcGIS~Online, Google~Maps~Platform,
    QGIS\,/\,OpenLayers) selon un format homogène.
  \item \textbf{Section~\ref{sec:bilan} — Bilan et Positionnement}~: tableau
    comparatif multicritère et positionnement de notre solution par rapport à l'existant.
\end{chapplan}

\lettrine[lines=2, loversize=0.05]{C}{e chapitre} pose les fondements théoriques
et contextuels nécessaires à la compréhension du travail réalisé.
La section~\ref{sec:concepts} présente les
concepts généraux mobilisés tout au long de ce rapport~: systèmes d'information
géographique, standards de services cartographiques, base de données spatiale,
architecture REST, contrôle d'accès basé sur les rôles et conteneurisation. La
section~\ref{sec:etat-art} dresse un état de l'art de quatre solutions existantes
répondant partiellement ou totalement aux besoins de l'EEC, analysées selon un
format homogène. La section~\ref{sec:bilan} synthétise cette analyse dans un
tableau comparatif et positionne notre solution par rapport à l'existant.

\section{Concepts Généraux}
\label{sec:concepts}

\subsection{Système d'Information Géographique (SIG)}

Un \kw{Système d'Information Géographique} (\kw{SIG}) est un système informatique
conçu pour la collecte, le stockage, la manipulation, l'analyse et la représentation
de données \kw{géoréférencées}, c'est-à-dire de données auxquelles est associée une
position sur la surface terrestre~\cite{longley2015}. Cette définition, établie
par Longley et al., est aujourd'hui l'une des plus citées dans la littérature
géomatique internationale. Un SIG intègre trois grandes dimensions~: la dimension
\textbf{spatiale} (coordonnées, géométries), la dimension \textbf{attributaire}
(données descriptives associées aux entités géographiques) et la dimension
\textbf{temporelle} (évolution des données dans le temps).

Dans le contexte de ce projet, le SIG constitue la couche centrale permettant de
lier les 553~paroisses importées, 137~districts et 22~régions synodales de l'EEC
à leurs coordonnées géographiques sur le territoire camerounais. Il permet non seulement
la visualisation cartographique de ces entités, mais aussi leur interrogation
spatiale~: identification des paroisses proches d'une localité, calcul de densité
pastorale par région, localisation des zones sous-dotées en ressources. Ces
opérations sont directement utiles au pilotage stratégique de l'institution.

\subsection{Standards OGC et Services Cartographiques Web}

L'\kw{Open Geospatial Consortium} (\kw{OGC}) est un organisme de normalisation
international fondé en 1994, regroupant des industriels, des organismes
gouvernementaux et des institutions académiques du domaine géospatial. Il
publie des \textbf{spécifications ouvertes} pour les services géographiques sur
le web, garantissant l'\textbf{interopérabilité} entre différents serveurs
cartographiques et clients web, indépendamment des fournisseurs et des
plateformes logicielles.

\subsubsection{Web Map Service (WMS)}

Le standard \textit{Web Map Service} (WMS) 1.3.0, publié par l'OGC en 2006,
définit une interface de service permettant à un client web de requêter des
représentations cartographiques sous forme d'images raster depuis un serveur
géographique~\cite{ogcwms2006}. Le client envoie une requête HTTP décrivant
l'étendue géographique souhaitée, les couches à afficher et les dimensions de
l'image~; le serveur retourne une image PNG ou JPEG directement superposable sur
n'importe quelle interface cartographique compatible. Dans ce projet, GeoServer
publie les couches des régions synodales, des paroisses et des œuvres sociales de
l'EEC en WMS~1.3.0, ce qui permet à l'interface Leaflet d'afficher ces couches
institutionnelles par-dessus le fond de carte OpenStreetMap.

\subsubsection{Web Feature Service (WFS)}

Le standard \textit{Web Feature Service} (WFS) 2.0.0, publié par l'OGC en 2010,
permet de requêter et de récupérer des données géographiques vectorielles ---
points, lignes, polygones --- sous forme structurée (GeoJSON ou GML) depuis un
serveur cartographique~\cite{ogcwfs2010}. Contrairement au WMS qui transmet des
images prérendues, le WFS retourne des données brutes directement manipulables
côté client. Ce projet exploite le WFS pour permettre à l'interface cartographique
de récupérer les géométries des régions synodales et d'appliquer des filtres
géographiques dynamiques selon le périmètre de l'utilisateur connecté.

\subsection{Base de Données Spatiale PostGIS}

\kw{PostGIS} est une extension spatiale open source pour le système de gestion de
base de données relationnelle \kw{PostgreSQL}~\cite{postgis2023}. Elle ajoute la
prise en charge de \textbf{types de données géographiques} conformes aux
spécifications OGC Simple Features --- géométries points (\texttt{PointField}),
lignes (\texttt{LineStringField}) et polygones (\texttt{PolygonField},
\texttt{MultiPolygonField}) --- ainsi que des \textbf{opérateurs spatiaux}
(intersections, calcul de distances, recouvrements) et des fonctions de
transformation de systèmes de coordonnées. PostGIS est le \textbf{standard de
facto} pour les bases de données spatiales open source, utilisé des systèmes
d'information institutionnels aux grandes applications de cartographie
collaborative.

Dans ce projet, PostGIS stocke la géométrie point de chaque paroisse et les
polygones des 22~régions synodales importés depuis le Shapefile source.
GeoDjango, la couche spatiale du framework Django, expose ces données via des
modèles Python et permet l'interrogation spatiale directement depuis l'ORM
(\textit{Object-Relational Mapper}), sans requêtes SQL manuelles.

\subsection{Architecture REST et Interface de Programmation (API)}

Le style architectural \kw{REST} (\textit{Representational State Transfer}) est
défini par Fielding comme un ensemble de contraintes architecturales pour les
systèmes distribués~\cite{fielding2000}~: séparation client-serveur, absence
d'état côté serveur (\textit{stateless}), mise en cache des réponses, interface
uniforme et système de couches. Une \kw{API REST} expose des ressources
identifiées par des URI, accessibles via les méthodes HTTP standard
(GET, POST, PUT, PATCH, DELETE), et retourne les données dans un format
standardisé, généralement \textbf{JSON}. Ce style garantit la \textbf{scalabilité},
la \textbf{maintenabilité} et l'\textbf{interopérabilité} des services web.

\kw{Django REST Framework} (\kw{DRF}) est la bibliothèque de référence pour
implémenter des API REST avec le framework Django~\cite{drf2024}. Dans ce projet, la couche
API du backend expose l'ensemble des ressources institutionnelles de l'EEC ---
paroisses, districts, régions, œuvres, ouvriers, statistiques, journaux d'audit
--- consommées par le frontend Next.js~14. Cette séparation stricte entre backend
API et frontend constitue un pilier de l'architecture à services découplés retenue.

\subsection{Contrôle d'Accès Basé sur les Rôles (RBAC)}

Le \kw{contrôle d'accès basé sur les rôles} (\kw{RBAC} --- \textit{Role-Based
Access Control}) est un modèle de sécurité dans lequel les droits d'accès aux
ressources d'un système sont accordés selon le \textbf{rôle} attribué à chaque
utilisateur, et non selon son identité individuelle. Un rôle représente un
ensemble cohérent de \textbf{permissions} correspondant à une fonction ou une
responsabilité dans l'organisation. Ce modèle réduit la complexité
administrative, minimise le risque d'\textbf{escalade de privilèges} et facilite
l'audit de sécurité. L'OWASP le recommande comme mécanisme d'autorisation de
référence pour les applications web~\cite{owasp2021}.

Dans ce projet, quatre rôles hiérarchiques alignés sur la structure synodale de
l'EEC ont été définis~:
\begin{itemize}[itemsep=2pt, topsep=4pt, leftmargin=1.5cm]
  \item \textbf{SUPER} (Bureau National)~: accès complet à l'ensemble des données
    de toutes les régions synodales~;
  \item \textbf{REGION} (Administrateur régional)~: accès limité aux entités de
    la région synodale dont l'utilisateur est responsable~;
  \item \textbf{DISTRICT} (Administrateur de district)~: accès limité aux entités
    du district ecclésial correspondant~;
  \item \textbf{PAROISSE} (Administrateur paroissial)~: accès en lecture étendue,
    modifications limitées à la paroisse de l'utilisateur.
\end{itemize}

\subsection{Conteneurisation --- Docker et Docker Compose}

La \kw{conteneurisation} est une technique de virtualisation légère qui isole une
application et l'ensemble de ses dépendances dans une unité standardisée appelée
\textbf{conteneur}. Contrairement à la virtualisation classique, un conteneur
partage le noyau du système d'exploitation hôte tout en disposant de son propre
système de fichiers, de ses processus et de sa pile réseau, le rendant bien plus
léger et plus rapide à démarrer qu'une machine virtuelle. \kw{Docker} est la
plateforme de conteneurisation la plus répandue, avec plusieurs millions de
déploiements en production à travers le monde~\cite{docker2024}.
\kw{Docker Compose} permet de définir et d'orchestrer des applications
multi-conteneurs dans un fichier de configuration déclaratif
(\texttt{docker-compose.yml}).

Ce projet déploie cinq services orchestrés par Docker Compose~:
\begin{enumerate}[itemsep=2pt, topsep=4pt, leftmargin=1.5cm]
  \item \texttt{backend} --- le serveur Django~5 avec GeoDjango et DRF~;
  \item \texttt{db} --- PostgreSQL~16 avec l'extension PostGIS~3.4~;
  \item \texttt{geoserver} --- GeoServer~2.26 pour la publication WMS/WFS~;
  \item \texttt{redis} --- le serveur de cache Redis~7~;
  \item \texttt{frontend} --- le serveur Next.js~14.
\end{enumerate}

Cette configuration garantit la reproductibilité de l'environnement d'exécution
sur n'importe quelle machine hôte et simplifie considérablement le déploiement
en production, un avantage décisif dans un contexte institutionnel à ressources
informatiques limitées.

% ----
\section{État de l'Art des Solutions Existantes}
\label{sec:etat-art}

Plusieurs catégories de solutions existent pour la gestion cartographique
institutionnelle. Cette section en analyse quatre représentatives~: la plateforme
GeoEEC développée antérieurement, deux plateformes cartographiques commerciales
majeures (ArcGIS Online et Google Maps Platform) et une solution libre de
référence (QGIS associé à OpenLayers). Pour chaque solution, le même format
d'analyse est appliqué~: description générale, illustration, fonctionnalités
principales, forces et faiblesses.

\subsection{Plateforme GeoEEC v1.0 (Version Antérieure)}

La plateforme GeoEEC v1.0 a été développée en 2025 dans le cadre d'un projet
académique antérieur à ce rapport. Elle constitue la référence la plus proche
de notre contexte puisqu'elle traite du même cas d'usage --- la cartographie des
entités de l'EEC --- et a été conçue avec une partie du même stack technologique.
Elle repose sur Django~5 (sans Docker), une authentification JWT
(\textit{JSON Web Token}), un frontend Next.js (Pages Router), une interface
cartographique Leaflet, et un chatbot intégré basé sur le modèle de langage
Ollama phi3:mini associé à LangChain et FAISS. Les données géographiques des
régions synodales sont stockées sous forme de fichiers GeoJSON statiques dans le
répertoire \texttt{public/} du frontend, sans serveur cartographique dédié.

\begin{figure}[H]
  \centering
  \fbox{\parbox{0.72\textwidth}{\centering\vspace{1.8cm}
    \textit{[Capture d'écran à intégrer~: Interface cartographique GeoEEC v1.0]}
    \vspace{1.8cm}}}
  \caption{Interface cartographique de la plateforme GeoEEC v1.0}
  \label{fig:geoeec-v1}
\end{figure}

\noindent\textbf{Fonctionnalités principales~:}
\begin{itemize}[itemsep=2pt, topsep=4pt, leftmargin=1.5cm]
  \item Carte interactive Leaflet affichant les paroisses, œuvres et ouvriers~;
  \item Import Excel des trois types d'entités avec gestion de la conversion des
    coordonnées sources (Coord\_x/Coord\_y)~;
  \item API REST CRUD pour les paroisses, œuvres et ouvriers~;
  \item Statistiques par région, district et type d'entité~; export PDF~;
  \item Chatbot IA contextuel (Ollama phi3:mini + LangChain + FAISS)~;
  \item Quatre rôles utilisateurs~: admin\_general, admin\_regional,
    utilisateur\_auth, visiteur.
\end{itemize}

\noindent\textbf{Forces~:}
\begin{itemize}[itemsep=2pt, topsep=4pt, leftmargin=1.5cm]
  \item Première implémentation fonctionnelle de la cartographie EEC, validant
    la faisabilité de l'approche Django + Leaflet~;
  \item Module d'import Excel opérationnel avec logique de dépivoter des données
    d'œuvres sociales~;
  \item Chatbot IA innovant permettant des requêtes en langage naturel sur les
    données institutionnelles.
\end{itemize}

\noindent\textbf{Faiblesses et limites~:}
\begin{itemize}[itemsep=2pt, topsep=4pt, leftmargin=1.5cm]
  \item GeoJSON statique~: les couches géographiques des régions synodales ne
    sont pas gérées par un serveur cartographique OGC~; tout changement de
    couche nécessite un redéploiement du frontend~;
  \item Authentification JWT~: les tokens transmis côté client sont plus
    vulnérables aux attaques XSS qu'une authentification par sessions
    \textit{HttpOnly}~;
  \item Absence de journal d'audit~: aucune traçabilité des modifications
    apportées aux données institutionnelles~;
  \item Déploiement non conteneurisé~: dépendances système non isolées,
    portabilité et maintenance complexes~;
  \item RBAC non aligné sur la hiérarchie synodale réelle~: les quatre rôles
    ne correspondent pas aux périmètres district et paroissial de l'EEC~;
  \item Absence de standards OGC~: interopérabilité cartographique limitée, pas
    de WMS ni de WFS exploitables depuis des clients tiers.
\end{itemize}

\subsection{ArcGIS Online (Esri)}

ArcGIS Online est la plateforme SIG en nuage de la société Esri, leader mondial
du logiciel géographique. Elle permet la création, la publication et l'analyse de
cartes interactives depuis un navigateur web, sans installation de logiciel, en
s'appuyant sur une infrastructure cloud hautement disponible. La plateforme intègre
des outils d'analyse spatiale avancée, de gestion de couches de données, de
création de tableaux de bord et de collaboration entre
équipes~\cite{arcgisonline2024}. Elle est largement déployée dans des organisations
gouvernementales, des ONG et des institutions académiques.

\begin{figure}[H]
  \centering
  \fbox{\parbox{0.72\textwidth}{\centering\vspace{1.8cm}
    \textit{[Capture d'écran à intégrer~: Interface ArcGIS Online]}
    \vspace{1.8cm}}}
  \caption{Plateforme cartographique ArcGIS Online (Esri, 2024)}
  \label{fig:arcgis-online}
\end{figure}

\noindent\textbf{Fonctionnalités principales~:}
\begin{itemize}[itemsep=2pt, topsep=4pt, leftmargin=1.5cm]
  \item Carte web interactive avec accès à de nombreux fonds de carte et couches
    de données mondiales~;
  \item Outils d'analyse spatiale avancés~: zones tampons, interpolations, analyses
    de réseaux, clustering~;
  \item Tableau de bord configurable avec indicateurs actualisables en temps réel~;
  \item Contrôle des permissions sur les couches (lecture, édition, partage)~;
  \item Applications mobiles intégrées pour la collecte de données terrain~;
  \item APIs REST compatibles avec les standards OGC WMS et WFS.
\end{itemize}

\noindent\textbf{Forces~:}
\begin{itemize}[itemsep=2pt, topsep=4pt, leftmargin=1.5cm]
  \item Plateforme mature avec un très large éventail de fonctionnalités SIG~;
  \item Interface utilisateur intuitive, accessible aux utilisateurs non
    techniciens~;
  \item Infrastructure cloud scalable avec haute disponibilité et support
    professionnel.
\end{itemize}

\noindent\textbf{Faiblesses et limites~:}
\begin{itemize}[itemsep=2pt, topsep=4pt, leftmargin=1.5cm]
  \item Licence commerciale onéreuse~: les abonnements annuels ArcGIS Online sont
    prohibitifs pour une institution confessionnelle à ressources limitées~;
  \item Aucune logique hiérarchique institutionnelle configurable~: pas de notion
    de périmètre synodal, de district ou de paroisse dans le modèle d'accès~;
  \item Données hébergées sur les serveurs d'Esri (États-Unis)~: dépendance à une
    infrastructure cloud étrangère, souveraineté des données compromise~;
  \item Absence de gestion intégrée des données métier EEC~: ouvriers
    ecclésiastiques, œuvres sociales et statistiques annuelles ne peuvent être
    gérés nativement.
\end{itemize}

\subsection{Google Maps Platform}

Google Maps Platform est l'ensemble des APIs cartographiques de Google, permettant
d'intégrer des cartes interactives, des fonctionnalités de géocodage et des
itinéraires dans des applications web et mobiles~\cite{googlemaps2024}. C'est la
solution cartographique web la plus utilisée mondialement. Elle s'appuie sur le
fond de carte Google, qui bénéficie d'une couverture très détaillée, y compris
pour les villes et localités camerounaises.

\begin{figure}[H]
  \centering
  \fbox{\parbox{0.72\textwidth}{\centering\vspace{1.8cm}
    \textit{[Capture d'écran à intégrer~: Google Maps Platform]}
    \vspace{1.8cm}}}
  \caption{Google Maps Platform --- Documentation et console développeur (Google LLC, 2024)}
  \label{fig:google-maps}
\end{figure}

\noindent\textbf{Fonctionnalités principales~:}
\begin{itemize}[itemsep=2pt, topsep=4pt, leftmargin=1.5cm]
  \item Cartographie interactive avec fond de carte Google haute résolution~;
  \item Géocodage (conversion adresse vers coordonnées GPS et inversement)~;
  \item Calcul d'itinéraires multimodaux~;
  \item Marqueurs personnalisés et couches de données~;
  \item APIs REST bien documentées, intégration simple pour les développeurs~;
  \item SDK mobiles Android et iOS.
\end{itemize}

\noindent\textbf{Forces~:}
\begin{itemize}[itemsep=2pt, topsep=4pt, leftmargin=1.5cm]
  \item Couverture cartographique mondiale très détaillée, performante pour le
    Cameroun~;
  \item API simple à intégrer, documentation exhaustive, large communauté de
    développeurs~;
  \item Géocodage très performant pour localiser des adresses camerounaises.
\end{itemize}

\noindent\textbf{Faiblesses et limites~:}
\begin{itemize}[itemsep=2pt, topsep=4pt, leftmargin=1.5cm]
  \item Modèle de facturation à l'usage~: au-delà du quota gratuit mensuel, les
    coûts peuvent devenir significatifs pour un usage institutionnel~;
  \item Pas de serveur cartographique OGC~: impossible de publier des couches
    propres en WMS ou WFS~;
  \item Aucune gestion des entités institutionnelles ni de la hiérarchie
    synodalo-ecclésiale~;
  \item Données cartographiques propriétaires hébergées sur les serveurs de
    Google~: dépendance totale et absence de souveraineté~;
  \item Pas de gestion des rôles, de journal d'audit ni d'import de données
    institutionnelles~;
  \item Conformité aux obligations légales locales et souveraineté des données
    problématiques pour une institution camerounaise.
\end{itemize}

\subsection{QGIS et OpenLayers}

QGIS est un Système d'Information Géographique libre et open source parmi les
plus utilisés au monde dans les secteurs académique, institutionnel et de la
coopération internationale~\cite{qgis2024}. Il offre des fonctionnalités
complètes d'édition, d'analyse et de visualisation cartographique depuis une
application bureau. OpenLayers est une bibliothèque JavaScript open source
permettant d'afficher des cartes interactives dans un navigateur web, avec support
natif des standards OGC WMS et WFS.

\begin{figure}[H]
  \centering
  \fbox{\parbox{0.72\textwidth}{\centering\vspace{1.8cm}
    \textit{[Capture d'écran à intégrer~: Interface QGIS 3.x]}
    \vspace{1.8cm}}}
  \caption{Interface du logiciel SIG bureau QGIS (version 3.x)}
  \label{fig:qgis}
\end{figure}

\noindent\textbf{Fonctionnalités principales~:}
\begin{itemize}[itemsep=2pt, topsep=4pt, leftmargin=1.5cm]
  \item Édition et analyse cartographique complète (QGIS)~;
  \item Symbologie avancée et mise en page de cartes imprimées~;
  \item Traitement géospatial avancé~: zones tampons, interpolations, analyses de
    réseau, calculs rasters~;
  \item Affichage de couches WMS et WFS dans le navigateur (OpenLayers)~;
  \item Support natif des formats SHP, GeoJSON, PostGIS et GeoTIFF.
\end{itemize}

\noindent\textbf{Forces~:}
\begin{itemize}[itemsep=2pt, topsep=4pt, leftmargin=1.5cm]
  \item Open source, sans coût de licence, avec une large communauté active et
    une documentation exhaustive~;
  \item Fonctionnalités SIG professionnelles complètes~;
  \item Haute interopérabilité grâce au support natif des standards OGC.
\end{itemize}

\noindent\textbf{Faiblesses et limites~:}
\begin{itemize}[itemsep=2pt, topsep=4pt, leftmargin=1.5cm]
  \item QGIS est une application bureau~: son déploiement en usage institutionnel
    web partagé est impraticable~;
  \item Aucun contrôle d'accès multi-utilisateurs basé sur des périmètres
    institutionnels hiérarchisés~;
  \item OpenLayers seul ne fournit ni backend, ni gestion de rôles, ni import de
    données institutionnelles, ni statistiques~;
  \item Nécessite des compétences SIG spécialisées~: inadapté aux administrateurs
    régionaux et paroissiaux non techniciens de l'EEC~;
  \item Déploiement complexe en environnement institutionnel à infrastructure
    informatique limitée.
\end{itemize}

% ----
\section{Bilan et Positionnement}
\label{sec:bilan}

\subsection{Tableau Comparatif des Solutions}

Le tableau~\ref{tab:comparatif} confronte les quatre solutions analysées aux onze
critères jugés déterminants pour une plateforme institutionnelle cartographique de
l'EEC. Les symboles employés sont les suivants~: \oui~critère entièrement satisfait,
\non~critère non satisfait, \partiel~critère partiellement satisfait, N/A~critère
sans objet pour cette catégorie de solution.

\begin{table}[H]
\centering
\caption{Tableau comparatif des solutions existantes et de la solution proposée}
\label{tab:comparatif}
\resizebox{\linewidth}{!}{%
\begin{tabular}{>{\bfseries}m{4.8cm}
                >{\centering\arraybackslash}m{2.0cm}
                >{\centering\arraybackslash}m{2.1cm}
                >{\centering\arraybackslash}m{2.1cm}
                >{\centering\arraybackslash}m{2.0cm}
                >{\centering\arraybackslash}m{2.2cm}}
\toprule
\textbf{Critère}
  & \textbf{GeoEEC v1.0}
  & \textbf{ArcGIS Online}
  & \textbf{Google Maps}
  & \textbf{QGIS / OL}
  & \textbf{Notre solution} \\
\midrule
Plateforme web multi-utilisateurs
  & \oui & \oui & \partiel & \non & \oui \\
\addlinespace[3pt]
Hiérarchie institutionnelle configurable
  & \partiel & \non & \non & \non & \oui \\
\addlinespace[3pt]
Authentification sécurisée (sessions)
  & \non & \oui & \oui & \non & \oui \\
\addlinespace[3pt]
RBAC à 4 niveaux synodaux
  & \partiel & \non & \non & \non & \oui \\
\addlinespace[3pt]
Serveur cartographique OGC (WMS/WFS)
  & \non & \oui & \non & \oui & \oui \\
\addlinespace[3pt]
Base de données spatiale (PostGIS)
  & \partiel & \oui & \oui & \oui & \oui \\
\addlinespace[3pt]
Journal d'audit des modifications
  & \non & \partiel & \non & \non & \oui \\
\addlinespace[3pt]
Import données Excel institutionnel
  & \oui & \non & \non & \non & \oui \\
\addlinespace[3pt]
Déploiement conteneurisé (Docker)
  & \non & N/A & N/A & \non & \oui \\
\addlinespace[3pt]
Logiciel libre / sans licence commerciale
  & \oui & \non & \non & \oui & \oui \\
\addlinespace[3pt]
Adapté au contexte institutionnel EEC
  & \partiel & \non & \non & \non & \oui \\
\bottomrule
\end{tabular}}
\end{table}

\subsection{Positionnement de Notre Solution}

L'analyse comparative met en évidence que notre solution est la seule à réunir
simultanément l'ensemble des caractéristiques requises pour une plateforme
institutionnelle cartographique au service de l'EEC.

Les plateformes commerciales --- ArcGIS Online et Google Maps Platform --- offrent
des capacités cartographiques puissantes mais ne répondent ni aux contraintes
économiques de l'institution ni à ses exigences de hiérarchie institutionnelle.
Leurs modèles de licences commerciales et leur dépendance à des infrastructures
cloud étrangères les rendent inadaptées à une institution confessionnelle opérant
au Cameroun. Elles ne proposent aucun mécanisme de gestion des données métier
spécifiques à l'EEC --- ouvriers ecclésiastiques, œuvres sociales, statistiques
annuelles des paroisses --- ni de contrôle d'accès structuré selon les périmètres
synodaux.

QGIS et OpenLayers, bien que techniquement solides et libres de droits, ne
constituent pas une plateforme web institutionnelle multi-utilisateurs~: QGIS est
une application bureau inadaptée aux administrateurs régionaux non techniciens, et
OpenLayers est une bibliothèque de rendu client sans backend, sans gestion de
données et sans contrôle d'accès.

La plateforme GeoEEC v1.0 constitue le précédent le plus proche de notre contexte,
mais présente des lacunes critiques identifiées lors de l'audit technique~:
vulnérabilité de l'authentification JWT par rapport aux sessions HttpOnly, absence
de journal d'audit, couches géographiques statiques non publiées via un serveur
OGC, RBAC non aligné sur la hiérarchie synodale réelle à quatre niveaux, et absence
de conteneurisation.

Notre solution apporte une valeur ajoutée précise et documentée sur chacun de ces
axes. Elle est la première à implémenter simultanément~: une architecture de services
découplés et entièrement conteneurisée, un serveur cartographique OGC natif
(GeoServer~2.26) publiant des couches WMS et WFS interopérables, une base de données
spatiale PostGIS avec gestion des géométries institutionnelles, une authentification
sécurisée par sessions Django (\textit{HttpOnly}, \textit{Secure},
\textit{SameSite=Strict}) et un RBAC exactement aligné sur les quatre niveaux
hiérarchiques de l'EEC, avec un journal d'audit complet et automatique de toutes
les opérations d'écriture sur les données.

\medskip
\noindent\textbf{Bilan du chapitre.}
Ce chapitre a établi les fondements théoriques du travail réalisé et mis en
évidence, par une analyse comparative rigoureuse, la pertinence et l'originalité
de l'approche proposée. Les six concepts fondamentaux --- SIG, standards OGC
(WMS/WFS), PostGIS, architecture REST, RBAC et Docker --- constituent le socle sur
lequel repose l'architecture décrite dans les chapitres suivants. L'état de l'art
de quatre solutions existantes a démontré qu'aucune d'elles ne couvre simultanément
l'ensemble des exigences institutionnelles de l'EEC. Le \textbf{Chapitre~2} procédera
à l'analyse détaillée des besoins du système à travers les exigences fonctionnelles
et non-fonctionnelles et les modélisations UML, puis présentera les décisions de
conception de l'architecture technique retenue.

\clearpage

% ============================================================
%  CHAPITRE 2 — ANALYSE ET CONCEPTION
% ============================================================
\chapter{Analyse et Conception du Système de Géolocalisation EEC}
\label{chap:analyse-conception}

\begin{chapplan}
  \item \textbf{Section~\ref{sec:analyse} --- Analyse des Besoins}~:
    \begin{itemize}[noitemsep, topsep=3pt, leftmargin=1.2cm]
      \item ingénierie des exigences~: 22~exigences fonctionnelles et 10~non-fonctionnelles~;
      \item diagramme de contexte~: six acteurs externes identifiés~;
      \item diagramme de packages~: trois couches (frontend, backend, infrastructure)~;
      \item cas d'utilisation UC01--UC06 avec tableaux de description structurés~;
      \item diagramme de classes métier~: huit entités du domaine institutionnel~;
      \item diagramme d'activité~: flux d'import de données depuis Excel.
    \end{itemize}
  \item \textbf{Section~\ref{sec:conception} --- Conception du Système}~:
    \begin{itemize}[noitemsep, topsep=3pt, leftmargin=1.2cm]
      \item architecture technique~: cinq services Docker découplés et leurs communications~;
      \item modèle de données~: quinze tables PostGIS avec champs géospatiaux~;
      \item diagramme de séquence~: flux d'authentification sécurisée par sessions Django~;
      \item diagramme de classes technique~: modèles Django et relations inter-classes.
    \end{itemize}
\end{chapplan}

\lettrine[lines=2, loversize=0.05]{L}{a} compréhension précise des besoins et
la définition rigoureuse d'une architecture sont les deux prérequis fondamentaux
d'une implémentation réussie. Ce chapitre traduit les contraintes institutionnelles
et techniques identifiées au Chapitre~1 en exigences mesurables et en décisions
de conception concrètes. La section~\ref{sec:analyse} formalise les besoins du
système à travers un tableau d'exigences fonctionnelles et non-fonctionnelles,
puis une série de modèles UML couvrant les acteurs, les modules, les cas
d'utilisation, les entités du domaine et les flux de traitement.
La section~\ref{sec:conception} s'appuie sur cette analyse pour définir
l'architecture technique retenue, le schéma de la base de données spatiale,
les diagrammes de séquence et de classes de l'implémentation.

%------------------------------------------------------------
\section{Analyse des Besoins}
\label{sec:analyse}
%------------------------------------------------------------

\subsection{Ingénierie des Exigences}
\label{ssec:exigences}

L'ingénierie des exigences consiste à identifier, formaliser et prioriser les
propriétés que le système doit satisfaire avant toute activité de conception.
Deux catégories sont distinguées~: les \textbf{exigences fonctionnelles}
(\textbf{EF}), qui décrivent ce que le système doit réaliser, et les
\textbf{exigences non-fonctionnelles} (\textbf{ENF}), qui portent sur les
contraintes qualitatives de fonctionnement (sécurité, performance, portabilité,
interopérabilité). Ces exigences sont directement dérivées des six objectifs
formulés en Introduction Générale et des lacunes identifiées dans l'état de
l'art du Chapitre~1.

\subsubsection{Exigences Fonctionnelles}

Le Tableau~\ref{tab:ef} recense les vingt-deux exigences fonctionnelles du système,
organisées par domaine~: authentification et contrôle d'accès, cartographie,
administration des entités, import/export, statistiques, traçabilité, gestion
des comptes et services du visiteur authentifié (EF19--EF22).

\setlength{\tabcolsep}{5pt}
\renewcommand{\arraystretch}{1.22}
\begin{longtable}{%
  >{\bfseries\color{eccdarkgreen}}m{1.2cm}
  m{7.8cm}
  m{3.5cm}
  >{\centering\arraybackslash\bfseries}m{1.7cm}}
\caption{Exigences fonctionnelles du système}\label{tab:ef} \\
\toprule
\rowcolor{eecgreen}
\multicolumn{1}{>{\bfseries\color{white}}m{1.2cm}}{ID} &
\multicolumn{1}{>{\bfseries\color{white}}m{7.8cm}}{Description} &
\multicolumn{1}{>{\bfseries\color{white}}m{3.5cm}}{Objectif} &
\multicolumn{1}{>{\bfseries\color{white}\centering\arraybackslash}m{1.7cm}}{Priorité} \\
\midrule
\endfirsthead
\caption[]{\textit{(suite)~Exigences fonctionnelles du système}} \\
\toprule
\rowcolor{eecgreen}
\multicolumn{1}{>{\bfseries\color{white}}m{1.2cm}}{ID} &
\multicolumn{1}{>{\bfseries\color{white}}m{7.8cm}}{Description} &
\multicolumn{1}{>{\bfseries\color{white}}m{3.5cm}}{Objectif} &
\multicolumn{1}{>{\bfseries\color{white}\centering\arraybackslash}m{1.7cm}}{Priorité} \\
\midrule
\endhead
\midrule
\multicolumn{4}{r}{\small\textit{(suite page suivante)}} \\
\endfoot
\bottomrule
\endlastfoot
\rowcolor{eeclightgreen!60}
EF01 & Le système doit permettre à tout utilisateur d'accéder à un formulaire
  de connexion sécurisé et de se déconnecter en invalidant la session active. &
  Sécurité & Haute \\
EF02 & Le système doit gérer cinq rôles hiérarchiques --- SUPER, REGION,
  DISTRICT, PAROISSE, VISITEUR --- chacun disposant d'un périmètre d'accès
  aux données strictement défini. & RBAC synodal & Haute \\
\rowcolor{eeclightgreen!60}
EF03 & Le système doit afficher une carte interactive Leaflet.js superposant
  les couches institutionnelles publiées par GeoServer~: paroisses (points),
  régions synodales (polygones) et œuvres (points). & Cartographie & Haute \\
EF04 & La carte doit regrouper automatiquement les marqueurs proches en clusters
  numériques, révélant les entités individuelles à mesure que l'utilisateur
  zoome. & Cartographie & Haute \\
\rowcolor{eeclightgreen!60}
EF05 & La carte doit proposer des filtres hiérarchiques multicritères~:
  sélection d'une région synodale, d'un district et d'un type d'entité
  (paroisse, œuvre, ouvrier). & Cartographie & Haute \\
EF06 & Un clic sur un marqueur de la carte doit ouvrir une popup affichant
  les informations détaillées de l'entité (nom, type, contacts, statut GPS). &
  Cartographie & Haute \\
\rowcolor{eeclightgreen!60}
EF07 & Les administrateurs autorisés doivent pouvoir créer, lire, modifier et
  supprimer des paroisses dans les limites de leur périmètre synodal. &
  Administration & Haute \\
EF08 & Les administrateurs autorisés doivent pouvoir créer, lire, modifier et
  supprimer des œuvres sociales et éducatives dans leur périmètre. &
  Administration & Haute \\
\rowcolor{eeclightgreen!60}
EF09 & Les administrateurs autorisés doivent pouvoir créer, lire, modifier et
  supprimer des ouvriers ecclésiastiques et gérer leurs affectations à une
  paroisse. & Administration & Haute \\
EF10 & Le système doit permettre l'import groupé des paroisses, ouvriers et
  œuvres depuis des fichiers Excel fournis par l'EEC, avec validation des
  données avant toute insertion en base de données. & Import & Haute \\
\rowcolor{eeclightgreen!60}
EF11 & Le système doit permettre l'export des listes d'entités au format
  Excel (.xlsx) et PDF depuis l'interface d'administration. & Export & Moyenne \\
EF12 & Le tableau de bord doit afficher des indicateurs synthétiques~:
  nombre de paroisses, districts, ouvriers, œuvres et taux de couverture
  GPS par région. & Statistiques & Haute \\
\rowcolor{eeclightgreen!60}
EF13 & Le système doit permettre la saisie, la modification et la
  consultation des statistiques annuelles de chaque paroisse (membres,
  baptisés, enfants, finances). & Statistiques & Haute \\
EF14 & Toute opération d'écriture (création, modification, suppression) sur
  les données institutionnelles doit être enregistrée automatiquement dans
  un journal d'audit horodaté. & Traçabilité & Haute \\
\rowcolor{eeclightgreen!60}
EF15 & Le journal d'audit doit être consultable par le rôle SUPER, avec des
  filtres par type d'action, entité concernée et période temporelle. &
  Traçabilité & Moyenne \\
EF16 & Le système doit proposer une recherche d'entités (paroisses, œuvres,
  ouvriers) par nom depuis la carte et depuis l'interface d'administration. &
  Ergonomie & Moyenne \\
\rowcolor{eeclightgreen!60}
EF17 & Les visiteurs non authentifiés doivent pouvoir consulter la carte
  publique et les informations générales des paroisses sans accéder au module
  d'administration. & Accès public & Moyenne \\
EF18 & L'administrateur SUPER doit pouvoir créer, désactiver et paramétrer
  les comptes utilisateurs (rôle, périmètre synodal, mot de passe). &
  Gestion comptes & Haute \\
\rowcolor{eeclightgreen!60}
EF19 & Le système doit permettre à un visiteur de s'inscrire avec un email
  et un mot de passe~; un profil visiteur est créé automatiquement à
  l'inscription et persiste entre les sessions. &
  Visiteur auth. & Haute \\
EF20 & Le visiteur authentifié doit pouvoir enregistrer des paroisses en
  favoris depuis la carte et les retrouver dans son espace personnel. &
  Visiteur auth. & Haute \\
\rowcolor{eeclightgreen!60}
EF21 & Le système doit calculer et afficher l'itinéraire vers la paroisse
  la plus proche du visiteur, avec distance (km) et durée estimée~; le calcul
  utilise la fonction PostGIS \texttt{Distance()}. &
  Visiteur auth. & Haute \\
EF22 & Le système doit conserver l'historique des recherches géographiques
  du visiteur authentifié et lui permettre de le consulter et de le vider
  depuis son profil personnel. &
  Visiteur auth. & Moyenne \\
\end{longtable}

\subsubsection{Exigences Non-Fonctionnelles}

Le Tableau~\ref{tab:enf} liste les dix exigences non-fonctionnelles définissant
les contraintes qualitatives de la plateforme sur les axes sécurité, performance,
interopérabilité, portabilité et maintenabilité.

\setlength{\tabcolsep}{5pt}
\renewcommand{\arraystretch}{1.22}
\begin{longtable}{%
  >{\bfseries\color{eccdarkgreen}}m{1.2cm}
  >{\bfseries}m{2.8cm}
  m{8.2cm}
  >{\centering\arraybackslash\bfseries}m{1.7cm}}
\caption{Exigences non-fonctionnelles du système}\label{tab:enf} \\
\toprule
\rowcolor{eecgreen}
\multicolumn{1}{>{\bfseries\color{white}}m{1.2cm}}{ID} &
\multicolumn{1}{>{\bfseries\color{white}}m{2.8cm}}{Domaine} &
\multicolumn{1}{>{\bfseries\color{white}}m{8.2cm}}{Description} &
\multicolumn{1}{>{\bfseries\color{white}\centering\arraybackslash}m{1.7cm}}{Priorité} \\
\midrule
\endfirsthead
\caption[]{\textit{(suite)~Exigences non-fonctionnelles du système}} \\
\toprule
\rowcolor{eecgreen}
\multicolumn{1}{>{\bfseries\color{white}}m{1.2cm}}{ID} &
\multicolumn{1}{>{\bfseries\color{white}}m{2.8cm}}{Domaine} &
\multicolumn{1}{>{\bfseries\color{white}}m{8.2cm}}{Description} &
\multicolumn{1}{>{\bfseries\color{white}\centering\arraybackslash}m{1.7cm}}{Priorité} \\
\midrule
\endhead
\midrule
\multicolumn{4}{r}{\small\textit{(suite page suivante)}} \\
\endfoot
\bottomrule
\endlastfoot
\rowcolor{eeclightgreen!60}
ENF01 & Sécurité & Les sessions sont gérées par des cookies Django
  \textit{HttpOnly}, \textit{Secure} et \textit{SameSite=Strict}. Chaque
  requête modifiante est protégée par un token CSRF. Les mots de passe sont
  hachés avec PBKDF2-SHA256. & Haute \\
ENF02 & Performance & Le temps de réponse des requêtes API standard (liste,
  détail) est inférieur à 2~secondes pour les jeux de données courants du
  système. & Haute \\
\rowcolor{eeclightgreen!60}
ENF03 & Disponibilité & Le déploiement Docker Compose sur un serveur
  institutionnel assure la continuité des services sans dépendance à une
  infrastructure cloud tierce. & Haute \\
ENF04 & Maintenabilité & L'architecture découplée (backend Django/DRF séparé
  du frontend Next.js) permet la mise à jour indépendante de chaque couche
  sans interruption des autres services. & Haute \\
\rowcolor{eeclightgreen!60}
ENF05 & Interopérabilité & Les couches géographiques sont publiées selon les
  standards OGC WMS~1.3.0 et WFS~2.0.0, garantissant leur consommation par
  tout client cartographique compatible. & Haute \\
ENF06 & Portabilité & La plateforme complète (cinq services) est conteneurisée
  avec Docker et Docker Compose, permettant le déploiement sur tout serveur
  Linux disposant de Docker. & Haute \\
\rowcolor{eeclightgreen!60}
ENF07 & Intégrité & L'import Excel déclenche une validation systématique des
  coordonnées GPS (plage territoire camerounais), du nommage des entités et
  des relations hiérarchiques avant toute insertion. & Haute \\
ENF08 & Traçabilité & Chaque entrée d'audit contient~: l'auteur, le type
  d'action (CREATE / UPDATE / DELETE), l'entité cible, l'identifiant de
  l'objet et l'horodatage UTC. & Haute \\
\rowcolor{eeclightgreen!60}
ENF09 & Ergonomie & L'interface d'administration est responsive et utilisable
  sur écrans desktop et tablette (résolution minimale~:
  1\,024~$\times$~768~px). & Moyenne \\
ENF10 & Extensibilité & Le schéma de données et l'API REST sont conçus pour
  accueillir de nouvelles régions synodales, types d'entités ou rôles sans
  refactorisation majeure du cœur du système. & Basse \\
\end{longtable}

%------------------------------------------------------------
\subsection{Diagramme de Contexte}
\label{ssec:contexte}

Le diagramme de contexte (Figure~\ref{fig:diag-contexte}) positionne la
plateforme EEC~Géolocalisation dans son environnement en l'abordant comme une
boîte noire~: il identifie les acteurs externes, la nature de leurs interactions
avec le système et les flux d'information échangés, sans détailler les composants
internes. Six acteurs externes sont identifiés.

\begin{itemize}[itemsep=3pt, topsep=5pt, leftmargin=1.5cm]
  \item \textbf{Administrateur SUPER} (Bureau National)~: acteur le plus
    privilégié, il dispose d'un accès complet à toutes les données --- gestion
    de toutes les entités institutionnelles, import/export, consultation du
    journal d'audit, gestion des comptes utilisateurs et statistiques de
    toutes les régions synodales~;
  \item \textbf{Administrateur REGION}~: accède aux mêmes fonctions
    d'administration que le SUPER, mais dans les limites strictes de sa
    région synodale~; il peut également consulter les analytiques régionales~;
  \item \textbf{Administrateurs DISTRICT et PAROISSE}~: accèdent aux fonctions
    CRUD dans le périmètre de leur district ou paroisse respectif~;
  \item \textbf{Visiteur Non Authentifié}~: accède uniquement à la carte
    publique et aux informations générales des paroisses, sans aucun accès
    au module d'administration ni aux services personnalisés~;
  \item \textbf{Visiteur Authentifié}~: visiteur inscrit qui bénéficie de
    services supplémentaires depuis la carte publique --- enregistrement de
    paroisses en favoris, calcul d'itinéraire vers la paroisse la plus proche
    (PostGIS \texttt{Distance()}), et historique de ses recherches~;
  \item \textbf{GeoServer} (système cartographique)~: système externe qui reçoit
    les données géographiques depuis le backend Django et répond aux requêtes
    WMS~1.3.0 et WFS~2.0.0 émises par le frontend Leaflet.
\end{itemize}

\begin{figure}[H]
  \centering
  \fbox{\parbox{0.88\textwidth}{\centering\vspace{2.8cm}
    \textit{[Diagramme de contexte --- à intégrer après génération draw.io]}%
    \vspace{2.8cm}}}
  \caption{Diagramme de contexte du système de géolocalisation EEC}
  \label{fig:diag-contexte}
\end{figure}

%------------------------------------------------------------
\subsection{Diagramme de Packages}
\label{ssec:packages}

Le diagramme de packages (Figure~\ref{fig:diag-packages}) décompose le système
en grands modules et représente leurs dépendances. Trois couches principales
sont identifiées.

\noindent\textbf{Couche frontend (Next.js~14)~:}
\begin{itemize}[itemsep=2pt, topsep=3pt, leftmargin=1.5cm]
  \item \texttt{app/(public)}~: interface publique avec carte Leaflet, clustering,
    filtres et panneau de détail~;
  \item \texttt{app/admin/(general)}~: tableau de bord SUPER/REGION (paroisses,
    districts, régions, œuvres, ouvriers, statistiques, audit, comptes,
    import/export)~;
  \item \texttt{app/admin/district}~: tableau de bord DISTRICT (accès restreint
    au périmètre du district)~;
  \item \texttt{app/admin/paroisse}~: tableau de bord PAROISSE (accès limité à
    la paroisse de l'utilisateur).
\end{itemize}

\noindent\textbf{Couche backend (Django~5 + DRF)~:}
\begin{itemize}[itemsep=2pt, topsep=3pt, leftmargin=1.5cm]
  \item \texttt{apps/paroisses}, \texttt{apps/districts}, \texttt{apps/regions}~:
    modèles géospatiaux, sérialiseurs et endpoints CRUD avec filtrage RBAC~;
  \item \texttt{apps/oeuvres}, \texttt{apps/ouvriers}~: gestion des entités
    sociales et du personnel ecclésiastique~;
  \item \texttt{apps/statistiques}~: saisie et consultation des statistiques
    annuelles par paroisse~;
  \item \texttt{apps/comptes}~: gestion des comptes utilisateurs, authentification
    par sessions Django et middleware RBAC~;
  \item \texttt{apps/journal}~: enregistrement automatique via signaux Django
    de toutes les opérations d'écriture~;
  \item \texttt{apps/visiteurs}~: profil visiteur public, historique de
    recherche et itinéraires personnels.
\end{itemize}

\noindent\textbf{Couche infrastructure~:} PostgreSQL~16 + PostGIS~3.4 (persistance
spatiale), GeoServer~2.26 (publication WMS/WFS), Redis~7 (cache), Docker Compose
(orchestration des cinq services).

\begin{figure}[H]
  \centering
  \fbox{\parbox{0.88\textwidth}{\centering\vspace{2.8cm}
    \textit{[Diagramme de packages --- à intégrer après génération draw.io]}%
    \vspace{2.8cm}}}
  \caption{Diagramme de packages du système de géolocalisation EEC}
  \label{fig:diag-packages}
\end{figure}

%------------------------------------------------------------
\subsection{Diagrammes de Cas d'Utilisation}
\label{ssec:uc}

Les diagrammes de cas d'utilisation modélisent les interactions entre les acteurs
et le système du point de vue fonctionnel. La Figure~\ref{fig:diag-uc} présente
le diagramme global consolidé de la plateforme, mettant en évidence les douze
cas d'utilisation principaux et leur distribution selon les six acteurs identifiés.
Les sous-sections suivantes détaillent les six cas d'utilisation les plus
représentatifs à travers des tableaux de description structurés~: UC01
(authentification), UC02 (carte interactive), UC03 (gestion des paroisses),
UC04 (import Excel), UC05 (journal d'audit) et UC06 (espace visiteur
authentifié).

\begin{figure}[H]
  \centering
  \fbox{\parbox{0.88\textwidth}{\centering\vspace{3.5cm}
    \textit{[Diagramme de cas d'utilisation global --- à intégrer après génération draw.io]}%
    \vspace{3.5cm}}}
  \caption{Diagramme global de cas d'utilisation de la plateforme EEC}
  \label{fig:diag-uc}
\end{figure}

\subsubsection{UC01 --- S'authentifier}

\begin{table}[H]
\centering
\caption{Description du cas d'utilisation UC01 --- S'authentifier}
\label{tab:uc01}
\small
\renewcommand{\arraystretch}{1.45}
\setlength{\tabcolsep}{6pt}
\begin{tabular}{>{\bfseries\color{eccdarkgreen}}p{3.0cm} p{10.8cm}}
\toprule
\rowcolor{eecgreen}
\multicolumn{2}{>{\bfseries\color{white}\centering\arraybackslash}p{13.8cm}}{UC01~--- S'authentifier} \\
\midrule
Acteurs & Tous les utilisateurs enregistrés~: SUPER, REGION, DISTRICT, PAROISSE, Visiteur Authentifié \\
\rowcolor{eeclightgreen!60}
Préconditions & Le compte de l'utilisateur existe et est actif dans le système \\
Scénario principal &
  1.~L'utilisateur navigue vers la page \texttt{/admin/login}.\newline
  2.~Il saisit son nom d'utilisateur et son mot de passe.\newline
  3.~Il soumet le formulaire de connexion.\newline
  4.~Le frontend envoie les credentials à l'API \texttt{/api/auth/login/}.\newline
  5.~Django vérifie les credentials, crée une session et transmet le cookie
     \texttt{sessionid} (\textit{HttpOnly, SameSite=Strict}).\newline
  6.~Le frontend redirige l'utilisateur vers le tableau de bord correspondant
     à son rôle. \\
\rowcolor{eeclightgreen!60}
Scénarios alternatifs &
  \textbf{Alt.~4a --- Credentials invalides}~: HTTP~401, message d'erreur
  affiché, utilisateur maintenu sur la page de connexion.\newline
  \textbf{Alt.~4b --- Compte désactivé}~: HTTP~403, message "Compte désactivé,
  contacter l'administrateur". \\
Postconditions & Session Django active, cookie \textit{HttpOnly} transmis,
  utilisateur redirigé vers son tableau de bord selon son rôle. \\
\bottomrule
\end{tabular}
\end{table}

\subsubsection{UC02 --- Consulter la Carte Interactive}

\begin{table}[H]
\centering
\caption{Description du cas d'utilisation UC02 --- Consulter la carte interactive}
\label{tab:uc02}
\small
\renewcommand{\arraystretch}{1.45}
\setlength{\tabcolsep}{6pt}
\begin{tabular}{>{\bfseries\color{eccdarkgreen}}p{3.0cm} p{10.8cm}}
\toprule
\rowcolor{eecgreen}
\multicolumn{2}{>{\bfseries\color{white}\centering\arraybackslash}p{13.8cm}}{UC02~--- Consulter la carte interactive} \\
\midrule
Acteurs & Visiteur Non Authentifié, Visiteur Authentifié et tous les rôles administratifs (SUPER, REGION, DISTRICT, PAROISSE) \\
\rowcolor{eeclightgreen!60}
Préconditions & Aucune (accès public pour le VISITEUR) \\
Scénario principal &
  1.~L'utilisateur accède à la page de la carte (\texttt{/} ou
     \texttt{/admin/map}).\newline
  2.~Le frontend Leaflet charge le fond de carte OpenStreetMap.\newline
  3.~Le frontend envoie des requêtes WMS à GeoServer pour les couches
     institutionnelles.\newline
  4.~Les marqueurs des paroisses et des œuvres s'affichent avec clustering
     automatique.\newline
  5.~L'utilisateur applique des filtres~: région synodale, district, type
     d'entité.\newline
  6.~La carte se met à jour dynamiquement selon les filtres sélectionnés.\newline
  7.~L'utilisateur clique sur un marqueur~; une popup s'ouvre avec les
     détails de l'entité. \\
\rowcolor{eeclightgreen!60}
Scénarios alternatifs &
  \textbf{Alt.~3a --- GeoServer inaccessible}~: les couches OGC ne s'affichent
  pas~; un avertissement est affiché~; le fond OpenStreetMap reste disponible.\newline
  \textbf{Alt.~6a --- Aucune entité pour les filtres}~: la carte reste vide
  avec un message informatif. \\
Postconditions & Les entités correspondant aux filtres actifs sont affichées
  sur la carte. \\
\bottomrule
\end{tabular}
\end{table}

\subsubsection{UC03 --- Gérer une Paroisse (CRUD)}

\begin{table}[H]
\centering
\caption{Description du cas d'utilisation UC03 --- Gérer une paroisse (CRUD)}
\label{tab:uc03}
\small
\renewcommand{\arraystretch}{1.45}
\setlength{\tabcolsep}{6pt}
\begin{tabular}{>{\bfseries\color{eccdarkgreen}}p{3.0cm} p{10.8cm}}
\toprule
\rowcolor{eecgreen}
\multicolumn{2}{>{\bfseries\color{white}\centering\arraybackslash}p{13.8cm}}{UC03~--- Gérer une paroisse (CRUD)} \\
\midrule
Acteurs & SUPER (toutes régions), REGION (sa région), DISTRICT (son district),
  PAROISSE (sa paroisse uniquement) \\
\rowcolor{eeclightgreen!60}
Préconditions & Utilisateur authentifié avec droits d'écriture sur la paroisse
  cible \\
Scénario principal &
  1.~L'administrateur accède à la liste des paroisses de son périmètre
     (\texttt{/admin/paroisses}).\newline
  2.~Il sélectionne l'action~: créer, modifier ou supprimer.\newline
  \textit{Création}~: il remplit le formulaire (nom, district, coordonnées GPS,
  statut) et soumet.\newline
  \textit{Modification}~: il modifie les champs voulus et soumet.\newline
  \textit{Suppression}~: il confirme l'action dans la boîte de dialogue.\newline
  3.~Le frontend envoie la requête (POST\,/\,PATCH\,/\,DELETE) vers
     \texttt{/api/paroisses/}.\newline
  4.~Le backend valide les données et vérifie les droits RBAC.\newline
  5.~L'opération est exécutée en base de données PostGIS.\newline
  6.~Un signal Django déclenche l'enregistrement dans le journal d'audit.\newline
  7.~Une réponse de succès est renvoyée~; la liste est actualisée. \\
\rowcolor{eeclightgreen!60}
Scénarios alternatifs &
  \textbf{Alt.~4a --- Violation RBAC}~: HTTP~403, message "Action non autorisée
  pour votre périmètre".\newline
  \textbf{Alt.~4b --- Données invalides}~: HTTP~400 avec les champs en erreur.\newline
  \textbf{Alt.~4c --- Coordonnées hors territoire}~: validation GPS échoue,
  message explicatif. \\
Postconditions & Paroisse créée/modifiée/supprimée en base~; action enregistrée
  dans le journal d'audit. \\
\bottomrule
\end{tabular}
\end{table}

\subsubsection{UC04 --- Importer des Données depuis Excel}

\begin{table}[H]
\centering
\caption{Description du cas d'utilisation UC04 --- Importer des données depuis Excel}
\label{tab:uc04}
\small
\renewcommand{\arraystretch}{1.45}
\setlength{\tabcolsep}{6pt}
\begin{tabular}{>{\bfseries\color{eccdarkgreen}}p{3.0cm} p{10.8cm}}
\toprule
\rowcolor{eecgreen}
\multicolumn{2}{>{\bfseries\color{white}\centering\arraybackslash}p{13.8cm}}{UC04~--- Importer des données depuis Excel} \\
\midrule
Acteurs & SUPER (import global), REGION (import pour sa région) \\
\rowcolor{eeclightgreen!60}
Préconditions & Utilisateur authentifié SUPER ou REGION~; fichier Excel fourni
  au format de l'EEC \\
Scénario principal &
  1.~L'administrateur accède au module d'import/export (\texttt{/admin/io}).\newline
  2.~Il sélectionne le type d'import~: paroisses, ouvriers ou œuvres.\newline
  3.~Il téléverse le fichier Excel (.xlsx) via le formulaire.\newline
  4.~Le backend parse le fichier avec \texttt{openpyxl} et valide chaque
     ligne~:\newline
  \quad --- vérification du format des colonnes (Coord\_x~= latitude,
     Coord\_y~= longitude)~;\newline
  \quad --- validation de la plage GPS (territoire camerounais)~;\newline
  \quad --- vérification de l'existence du district et de la région référencés~;\newline
  \quad --- détection des doublons (même nom + même district).\newline
  5.~Les lignes valides sont insérées en base de données~; les invalides sont
     consignées.\newline
  6.~Un rapport d'import est retourné~: lignes insérées, ignorées, erreurs
     détaillées par ligne. \\
\rowcolor{eeclightgreen!60}
Scénarios alternatifs &
  \textbf{Alt.~3a --- Format de fichier incorrect}~: HTTP~400, message "Format
  non reconnu, utilisez le modèle fourni".\newline
  \textbf{Alt.~5a --- Toutes les lignes invalides}~: aucune insertion~; rapport
  d'erreurs complet retourné.\newline
  \textbf{Alt.~5b --- Import partiel}~: les lignes valides sont insérées~; le
  rapport liste les lignes rejetées avec raison. \\
Postconditions & Entités valides enregistrées en base PostGIS~; entrée d'audit
  synthétique créée~; rapport de résultats affiché. \\
\bottomrule
\end{tabular}
\end{table}

\subsubsection{UC05 --- Consulter le Journal d'Audit}

\begin{table}[H]
\centering
\caption{Description du cas d'utilisation UC05 --- Consulter le journal d'audit}
\label{tab:uc05}
\small
\renewcommand{\arraystretch}{1.45}
\setlength{\tabcolsep}{6pt}
\begin{tabular}{>{\bfseries\color{eccdarkgreen}}p{3.0cm} p{10.8cm}}
\toprule
\rowcolor{eecgreen}
\multicolumn{2}{>{\bfseries\color{white}\centering\arraybackslash}p{13.8cm}}{UC05~--- Consulter le journal d'audit} \\
\midrule
Acteurs & SUPER uniquement \\
\rowcolor{eeclightgreen!60}
Préconditions & Utilisateur authentifié avec le rôle SUPER \\
Scénario principal &
  1.~L'administrateur SUPER navigue vers \texttt{/admin/journal}.\newline
  2.~La liste des entrées s'affiche~: auteur, type d'action, entité,
     horodatage UTC.\newline
  3.~Il applique des filtres~: type d'action (CREATE/UPDATE/DELETE), type
     d'entité, période.\newline
  4.~La liste filtrée se met à jour dynamiquement.\newline
  5.~Il clique sur une entrée pour consulter le détail~: valeurs avant et
     après la modification. \\
\rowcolor{eeclightgreen!60}
Scénarios alternatifs &
  \textbf{Alt.~1a --- Accès depuis un rôle non SUPER}~: HTTP~403, redirection
  vers le tableau de bord de l'utilisateur. \\
Postconditions & Les entrées d'audit correspondant aux filtres sont affichées~;
  aucune modification n'est apportée aux données. \\
\bottomrule
\end{tabular}
\end{table}

\subsubsection{UC06 --- S'inscrire et Utiliser l'Espace Visiteur Authentifié}

\begin{table}[htbp]
\centering
\caption{Description du cas d'utilisation UC06 --- Espace visiteur authentifié}
\label{tab:uc06}
\small
\renewcommand{\arraystretch}{1.45}
\setlength{\tabcolsep}{6pt}
\begin{tabular}{>{\bfseries\color{eccdarkgreen}}p{3.0cm} p{10.8cm}}
\toprule
\rowcolor{eecgreen}
\multicolumn{2}{>{\bfseries\color{white}\centering\arraybackslash}p{13.8cm}}{UC06~--- S'inscrire et utiliser l'espace visiteur authentifié} \\
\midrule
Acteurs & Visiteur Authentifié (inscrit via email/mot de passe) \\
\rowcolor{eeclightgreen!60}
Préconditions & Accès à la carte publique~; adresse email non encore enregistrée \\
Scénario principal &
  1.~Le visiteur clique sur «~S'inscrire~» depuis la carte publique.\newline
  2.~Il saisit son email et son mot de passe dans le formulaire d'inscription.\newline
  3.~Le frontend envoie les données à \texttt{/api/visiteurs/register/}.\newline
  4.~Django crée le compte et génère automatiquement un
     \textit{ProfilVisiteur} associé.\newline
  5.~Le visiteur est connecté~; son profil personnel devient accessible
     depuis la carte.\newline
  6.~Sur la carte, il clique sur «~Paroisse la plus proche~»~: le backend
     interroge PostGIS (\texttt{Distance()}) et retourne la paroisse la
     plus proche avec distance (km) et durée estimée ($\approx$~60~km/h).\newline
  7.~Il clique sur «~Enregistrer~» sur la fiche d'une paroisse~;
     l'enregistrement est sauvegardé dans \textit{ParoisseEnregistree}
     (\texttt{/api/visiteurs/paroisses/}).\newline
  8.~Il consulte son espace personnel~: liste des paroisses enregistrées,
     historique de ses recherches dans \textit{RechercheHistorique}. \\
\rowcolor{eeclightgreen!60}
Scénarios alternatifs &
  \textbf{Alt.~3a --- Email déjà utilisé}~: HTTP~400, message «~Un compte
  existe déjà avec cet email~».\newline
  \textbf{Alt.~6a --- Localisation non partagée}~: message d'invitation à
  activer la géolocalisation dans le navigateur~; fonctionnalité non disponible
  sans position GPS.\newline
  \textbf{Alt.~8a --- Vider l'historique}~: le visiteur supprime toutes les
  entrées de \textit{RechercheHistorique} via DELETE \texttt{/api/visiteurs/historique/}. \\
Postconditions & ProfilVisiteur créé~; paroisses enregistrées et historique
  persistés en base~; session visiteur active. \\
\bottomrule
\end{tabular}
\end{table}

%------------------------------------------------------------
\subsection{Diagramme de Classes Métier}
\label{ssec:classes-metier}

Le diagramme de classes métier (Figure~\ref{fig:diag-classes-metier}) modélise
les entités du domaine institutionnel de l'EEC et leurs relations, indépendamment
de tout détail d'implémentation technique. Huit entités principales sont
identifiées~:

\begin{itemize}[itemsep=2pt, topsep=4pt, leftmargin=1.5cm]
  \item \textbf{RegionSynodale}~: entité géographique de premier niveau, définie
    par un nom, un siège et un polygone géospatial~;
  \item \textbf{District}~: entité de second niveau, rattachée à une région
    synodale~;
  \item \textbf{Paroisse}~: entité de base de la présence de l'EEC, rattachée à
    un district, localisée par des coordonnées GPS~;
  \item \textbf{OEuvre}~: établissement social ou éducatif (école, centre
    médical, exploitation agropastorale, immeuble), rattaché à une paroisse~;
  \item \textbf{Ouvrier}~: personnel ecclésiastique (pasteur, évangéliste,
    diacre), affecté à une paroisse~;
  \item \textbf{StatistiqueAnnuelle}~: relevé chiffré annuel d'une paroisse
    (membres, baptisés, enfants, indicateurs financiers)~;
  \item \textbf{CompteUtilisateur}~: compte d'accès à la plateforme, associé
    à un rôle et à un périmètre synodal (région, district ou paroisse)~;
  \item \textbf{EntreeAudit}~: enregistrement immuable d'une opération
    d'écriture, horodaté et lié à l'auteur.
\end{itemize}

\begin{figure}[H]
  \centering
  \fbox{\parbox{0.88\textwidth}{\centering\vspace{3.0cm}
    \textit{[Diagramme de classes métier --- à intégrer après génération draw.io]}%
    \vspace{3.0cm}}}
  \caption{Diagramme de classes métier du domaine institutionnel EEC}
  \label{fig:diag-classes-metier}
\end{figure}

%------------------------------------------------------------
\subsection{Diagramme d'Activité}
\label{ssec:activite}

Le diagramme d'activité (Figure~\ref{fig:diag-activite}) modélise le flux de
traitement du processus d'import de données depuis Excel --- le processus métier
le plus critique et le plus complexe du système, car il constitue le seul point
d'entrée de masse pour les données institutionnelles et implique une chaîne de
validation multi-étapes avant toute persistance.

Le flux se déroule en cinq phases~: (1)~téléversement et validation du format
du fichier Excel~; (2)~extraction et validation ligne par ligne (coordonnées GPS,
cohérence hiérarchique, doublons)~; (3)~insertion des lignes valides en base de
données PostGIS~; (4)~enregistrement d'une entrée d'audit synthétique~;
(5)~génération et affichage du rapport d'import détaillé.

\begin{figure}[H]
  \centering
  \fbox{\parbox{0.88\textwidth}{\centering\vspace{3.0cm}
    \textit{[Diagramme d'activité --- Import Excel --- à intégrer après génération draw.io]}%
    \vspace{3.0cm}}}
  \caption{Diagramme d'activité du processus d'import de données depuis Excel}
  \label{fig:diag-activite}
\end{figure}

%============================================================
\section{Conception du Système}
\label{sec:conception}
%============================================================

\subsection{Architecture Technique}
\label{ssec:archi}

La Figure~\ref{fig:diag-archi} présente l'architecture technique retenue~: une
architecture à cinq services découplés, orchestrés par Docker Compose sur un
serveur unique. Chaque service est isolé dans son propre conteneur et communique
avec les autres via le réseau Docker interne \texttt{eec\_network}.

\noindent\textbf{Services et rôles~:}
\begin{itemize}[itemsep=2pt, topsep=4pt, leftmargin=1.5cm]
  \item \textbf{Frontend Next.js~16} (port~3004)~: application React côté
    serveur avec App Router, TypeScript et Tailwind CSS. Sert la carte publique
    Leaflet et les quatre interfaces d'administration. Communique avec le backend
    exclusivement via l'API REST (\texttt{/api/}), sans accès direct à la base
    de données~;
  \item \textbf{Backend Django~5 + DRF} (port~8000)~: API REST de référence.
    Expose 40+ endpoints (geo, auth, exports, imports, visiteurs, analytics,
    audit). Applique le RBAC à chaque requête via un middleware dédié. Gère les
    sessions dans Redis. Exécute les requêtes géospatiales via GeoDjango et
    PostGIS~;
  \item \textbf{PostgreSQL~16 + PostGIS~3.4} (port~5432)~: base de données
    relationnelle et spatiale. Stocke l'intégralité des données institutionnelles
    (paroisses avec \texttt{PointField}, régions avec \texttt{MultiPolygonField})
    et les données métier (ouvriers, œuvres, statistiques, audit)~;
  \item \textbf{GeoServer~2.26} (port~8080)~: serveur cartographique OGC. Publie
    les couches géographiques de la base PostGIS en WMS~1.3.0 (images raster)
    et WFS~2.0.0 (features vectorielles). Consommé directement par Leaflet dans
    le frontend~;
  \item \textbf{Redis~7} (port~6379)~: stockage en mémoire utilisé pour la
    gestion des sessions Django (\texttt{SESSION\_ENGINE =
    django.contrib.sessions.backends.cache}) et la mise en cache des requêtes
    fréquentes.
\end{itemize}

\noindent\textbf{Flux de communication~:} Le frontend Next.js adresse ses
requêtes API au backend Django. Le backend interroge PostgreSQL/PostGIS pour
toutes les opérations de données et Redis pour les sessions. GeoServer accède
directement à PostGIS via JDBC pour lire les couches géographiques et les
servir au frontend Leaflet. Aucun service n'est exposé directement à Internet
--- un reverse proxy (Nginx) centralise les accès externes en production.

\begin{figure}[htbp]
  \centering
  \fbox{\parbox{0.88\textwidth}{\centering\vspace{3.0cm}
    \textit{[Diagramme d'architecture technique --- cinq services Docker
    --- à intégrer après génération draw.io]}%
    \vspace{3.0cm}}}
  \caption{Architecture technique de la plateforme EEC Géolocalisation}
  \label{fig:diag-archi}
\end{figure}

%------------------------------------------------------------
\subsection{Modèle de Données}
\label{ssec:modele-bd}

Le modèle de données repose sur quinze tables PostGIS réparties en cinq
domaines fonctionnels. Le Tableau~\ref{tab:modele-bd} synthétise les tables
principales, leurs champs géospatiaux et leurs relations clés~; la
Figure~\ref{fig:diag-modele-bd} en présente le schéma entité-relation.

\setlength{\tabcolsep}{5pt}
\renewcommand{\arraystretch}{1.22}
\begin{longtable}{%
  >{\bfseries\color{eccdarkgreen}}m{3.5cm}
  m{4.5cm}
  m{6.3cm}}
\caption{Principales tables du modèle de données PostGIS}\label{tab:modele-bd} \\
\toprule
\rowcolor{eecgreen}
\multicolumn{1}{>{\bfseries\color{white}}m{3.5cm}}{Table} &
\multicolumn{1}{>{\bfseries\color{white}}m{4.5cm}}{Champ géospatial / clé} &
\multicolumn{1}{>{\bfseries\color{white}}m{6.3cm}}{Rôle et relations principales} \\
\midrule
\endfirsthead
\caption[]{\textit{(suite)~Modèle de données PostGIS}} \\
\toprule
\rowcolor{eecgreen}
\multicolumn{1}{>{\bfseries\color{white}}m{3.5cm}}{Table} &
\multicolumn{1}{>{\bfseries\color{white}}m{4.5cm}}{Champ géospatial / clé} &
\multicolumn{1}{>{\bfseries\color{white}}m{6.3cm}}{Rôle et relations principales} \\
\midrule
\endhead
\midrule
\multicolumn{3}{r}{\small\textit{(suite page suivante)}} \\
\endfoot
\bottomrule
\endlastfoot
\rowcolor{eeclightgreen!60}
geo\_regionsynodale & \texttt{polygone} MultiPolygonField (SRID=4326) & Entité de niveau 1~; racine de la hiérarchie synodale \\
geo\_district & FK région (PROTECT) & Entité de niveau 2~; 137 districts \\
\rowcolor{eeclightgreen!60}
geo\_paroisse & \texttt{localisation} PointField (SRID=4326) & Entité de base~; FK district (PROTECT)~; 565 paroisses \\
oeuvres\_oeuvre & \texttt{localisation} PointField (SRID=4326) & Établissements sociaux/éducatifs~; FK paroisse (SET\_NULL) \\
\rowcolor{eeclightgreen!60}
ouvriers\_ouvrier & --- & Personnel ecclésiastique~; FK paroisse (SET\_NULL) \\
statistiques\_stat & --- & Relevé annuel par paroisse~; FK paroisse (CASCADE) \\
\rowcolor{eeclightgreen!60}
accounts\_utilisateur & --- & Étend AbstractUser~; champs role, region, district, paroisse \\
accounts\_profil & --- & OneToOne vers Utilisateur~; informations de contact étendues \\
\rowcolor{eeclightgreen!60}
audit\_entreeaudit & --- & Journal immuable~; FK utilisateur (SET\_NULL) \\
visitors\_profilvisiteur & --- & OneToOne vers User~; créé auto à l'inscription \\
\rowcolor{eeclightgreen!60}
visitors\_paroisseenreg & --- & M2M Visiteur × Paroisse~; date d'enregistrement \\
visitors\_itineraire & \texttt{point\_depart} PointField & FK ProfilVisiteur (CASCADE)~; paroisse cible FK \\
\rowcolor{eeclightgreen!60}
visitors\_recherche & --- & Historique de recherche~; FK ProfilVisiteur (CASCADE) \\
exports\_rapport & --- & Métadonnées des exports générés (PDF, XLSX) \\
\rowcolor{eeclightgreen!60}
geo\_paroissevue & --- & Compteur de vues par paroisse~; FK paroisse (CASCADE) \\
\end{longtable}

\begin{figure}[htbp]
  \centering
  \fbox{\parbox{0.88\textwidth}{\centering\vspace{3.0cm}
    \textit{[Schéma entité-relation (ERD) --- à intégrer après génération draw.io]}%
    \vspace{3.0cm}}}
  \caption{Modèle de données --- schéma entité-relation de la plateforme EEC}
  \label{fig:diag-modele-bd}
\end{figure}

%------------------------------------------------------------
\subsection{Diagramme de Séquence}
\label{ssec:sequence}

Le diagramme de séquence (Figure~\ref{fig:diag-sequence}) modélise le flux
d'authentification sécurisée, scénario critique qui conditionne l'accès à
l'ensemble des fonctionnalités protégées. Ce flux illustre les interactions
entre le navigateur, le frontend Next.js, le middleware Django et la base
de données Redis.

Le flux se déroule en dix étapes~:
\begin{enumerate}[itemsep=2pt, topsep=4pt, leftmargin=1.8cm]
  \item Le navigateur émet \texttt{GET /admin/login} pour obtenir la page
    de connexion~;
  \item Next.js retourne la page et un token CSRF préchargé~;
  \item L'utilisateur saisit ses identifiants et soumet le formulaire~;
  \item Next.js émet \texttt{POST /api/auth/login/} avec les credentials
    et le token CSRF dans l'en-tête~;
  \item Le middleware Django valide le token CSRF~; en cas d'échec, HTTP~403~;
  \item Django interroge la table \texttt{accounts\_utilisateur} et vérifie
    le hash PBKDF2-SHA256 du mot de passe~;
  \item Si les credentials sont corrects, Django crée une session dans Redis
    (\texttt{SESSION\_ENGINE = django.contrib.sessions.backends.cache})~;
  \item Django retourne HTTP~200 avec le cookie \texttt{sessionid}
    (\textit{HttpOnly, SameSite=Strict})~;
  \item Next.js stocke le cookie (géré automatiquement par le navigateur~;
    jamais exposé à JavaScript)~;
  \item Next.js redirige vers le tableau de bord correspondant au rôle de
    l'utilisateur (\texttt{/admin/general}, \texttt{/admin/regional},
    \texttt{/admin/district} ou \texttt{/admin/paroisse}).
\end{enumerate}

\begin{figure}[htbp]
  \centering
  \fbox{\parbox{0.88\textwidth}{\centering\vspace{3.0cm}
    \textit{[Diagramme de séquence --- Authentification sécurisée ---
    à intégrer après génération draw.io]}%
    \vspace{3.0cm}}}
  \caption{Diagramme de séquence du flux d'authentification sécurisée}
  \label{fig:diag-sequence}
\end{figure}

%------------------------------------------------------------
\subsection{Diagramme de Classes Technique}
\label{ssec:classes-tech}

Le diagramme de classes technique (Figure~\ref{fig:diag-classes-tech})
représente l'implémentation concrète des entités sous forme de modèles Django,
avec leurs attributs, types Python/PostGIS et relations inter-classes.

\noindent\textbf{Modèles principaux et attributs~:}
\begin{itemize}[itemsep=2pt, topsep=4pt, leftmargin=1.5cm]
  \item \textbf{Paroisse} (app \texttt{geo})~: \texttt{nom} (CharField),
    \texttt{district} (FK → District, PROTECT), \texttt{localisation}
    (PointField, SRID=4326, null=True), \texttt{statut} (CharField, choix
    ACTIVE/INACTIVE), \texttt{has\_gps} (BooleanField)~;
  \item \textbf{RegionSynodale} (app \texttt{geo})~: \texttt{nom} (CharField),
    \texttt{siege} (CharField), \texttt{polygone} (MultiPolygonField,
    SRID=4326)~;
  \item \textbf{Utilisateur} (app \texttt{accounts})~: étend
    \texttt{AbstractUser}~; \texttt{role} (CharField, choix SUPER/REGION/
    DISTRICT/PAROISSE/VISITEUR), \texttt{region} (FK → RegionSynodale,
    SET\_NULL), \texttt{district} (FK → District, SET\_NULL), \texttt{paroisse}
    (FK → Paroisse, SET\_NULL)~;
  \item \textbf{ProfilVisiteur} (app \texttt{visitors})~: \texttt{user}
    (OneToOneField → Utilisateur, CASCADE), \texttt{date\_inscription}
    (DateTimeField, auto), \texttt{derniere\_position} (PointField, null=True)~;
  \item \textbf{ItinerairePersonnel} (app \texttt{visitors})~: \texttt{profil}
    (FK → ProfilVisiteur, CASCADE), \texttt{paroisse\_cible} (FK → Paroisse,
    CASCADE), \texttt{point\_depart} (PointField), \texttt{distance\_km}
    (FloatField), \texttt{duree\_min} (IntegerField)~;
  \item \textbf{EntreeAudit} (app \texttt{audit})~: \texttt{auteur} (FK →
    Utilisateur, SET\_NULL), \texttt{action} (CharField, choix CREATE/UPDATE/
    DELETE), \texttt{modele} (CharField), \texttt{objet\_id} (IntegerField),
    \texttt{horodatage} (DateTimeField, auto\_now\_add).
\end{itemize}

\noindent\textbf{Relations inter-classes~:} la relation Paroisse → District →
RegionSynodale forme la colonne vertébrale hiérarchique (PROTECT empêche la
suppression en cascade). Les relations Oeuvre et Ouvrier vers Paroisse sont
SET\_NULL (un ouvrier peut exister sans paroisse). EntreeAudit est reliée à
Utilisateur en SET\_NULL pour conserver les logs même après suppression du
compte. ProfilVisiteur est en OneToOne avec Utilisateur (CASCADE).

\begin{figure}[htbp]
  \centering
  \fbox{\parbox{0.88\textwidth}{\centering\vspace{3.0cm}
    \textit{[Diagramme de classes technique --- modèles Django ---
    à intégrer après génération draw.io]}%
    \vspace{3.0cm}}}
  \caption{Diagramme de classes technique des modèles Django de la plateforme}
  \label{fig:diag-classes-tech}
\end{figure}

\medskip
\noindent\textbf{Bilan du chapitre.}
Ce chapitre a formalisé l'ensemble des besoins de la plateforme en
vingt-deux exigences fonctionnelles et dix exigences non-fonctionnelles,
couvrant les domaines de l'authentification, de la cartographie, de
l'administration des entités, de l'import/export, des statistiques, de la
traçabilité et des services du visiteur authentifié. Les six modèles UML
d'analyse --- contexte (six acteurs), packages, cas d'utilisation
(UC01--UC06), classes métier et activité --- ont défini avec précision les
acteurs, les modules, les interactions et les entités du domaine institutionnel.
La section de conception a complété cette analyse en présentant l'architecture
technique à cinq services Docker, le modèle de données PostGIS (quinze tables),
le diagramme de séquence d'authentification sécurisée et le diagramme de
classes technique des modèles Django.
Le \textbf{Chapitre~3} présentera l'environnement de développement retenu,
justifiera les choix technologiques et exposera les résultats concrets de
l'implémentation.

\clearpage

% ============================================================
%  CHAPITRE 3 — IMPLÉMENTATION ET RÉSULTATS (À VENIR)
%  CONCLUSION GÉNÉRALE (À VENIR)
% ============================================================

% ============================================================
%  RÉFÉRENCES BIBLIOGRAPHIQUES
% ============================================================
\clearpage
\printbibliography[heading=bibintoc, title={Références Bibliographiques}]

\end{document}
